% ============================================================
%  my_preamble.sty
%  upLaTeX + dvipdfmx + Beamer 用プリアンブル
%
%  使い方（\documentclass の直後に）:
%    % ============================================================
%  my_preamble.sty
%  upLaTeX + dvipdfmx + Beamer 用プリアンブル
%
%  使い方（\documentclass の直後に）:
%    % ============================================================
%  my_preamble.sty
%  upLaTeX + dvipdfmx + Beamer 用プリアンブル
%
%  使い方（\documentclass の直後に）:
%    % ============================================================
%  my_preamble.sty
%  upLaTeX + dvipdfmx + Beamer 用プリアンブル
%
%  使い方（\documentclass の直後に）:
%    \input{my_preamble}
%
%  前提: \documentclass[dvipdfmx, ...]{beamer}
% ============================================================

% === 数学・記号 ==============================================
% amsmath / amsfonts は Beamer が内部でロード済みなので不要
% amssymb のみ Beamer が読まないため明示的にロードする
\usepackage{amssymb}

% === 図版・作図 ==============================================
% standalone: \input で standalone クラスのファイルを取り込む際に
%   プリアンブル（\documentclass 〜 \begin{document}）を自動的に無視し、
%   本文（TikZ コード等）だけを展開する。
%   → 各 TikZ ファイルを単体コンパイルしながら Beamer にも \input できる。
\usepackage{standalone}

\usepackage{tikz}
\usetikzlibrary{arrows.meta}               % Stealth 矢印等
\usetikzlibrary{decorations.pathreplacing} % 波括弧（DiD図等）
\usetikzlibrary{shapes.geometric, positioning}          % 楕円ノード等（因果推論図）

% PGFPlots: 関数プロット・データ可視化
\usepackage{pgfplots}
\pgfplotsset{compat=1.18}
\usepgfplotslibrary{groupplots}       % 複数パネル（上下・左右配置）
\usepgfplotslibrary{fillbetween}      % 2曲線間の塗りつぶし

% === TikZ 図共通カラー ========================================
% standalone プリアンブルは \input 時に無視されるため、
% 各 TikZ ファイルで使う色はここで定義する。
\definecolor{cfcol}  {RGB}{232, 232, 232}     % 観測不可（薄灰）
\definecolor{cfgray}  {RGB}{100, 140, 200}    % 反事実（薄い青）
\definecolor{expgreen} {RGB}{200, 230, 200}   % 拡張局面シェード
\definecolor{fillcol}{RGB}{180, 215, 255}     % 政策効果の差分領域
\definecolor{gdpblue}  {RGB}{30,  90,  180}   % GDP 線・AD 曲線
\definecolor{goldcol}  {RGB}{190, 140,  10}
\definecolor{recred}   {RGB}{240, 200, 200}   % 後退局面シェード
\definecolor{srasgreen}{RGB}{30,  140, 60}    % SRAS（短期総供給曲線）
\definecolor{trendred} {RGB}{190, 40,  40}    % トレンド線・LRAS
\definecolor{obscol} {RGB}{210, 228, 255}     % 観測可能（薄青）
\definecolor{watercol}{RGB}{175, 210, 240}

% graphicx は Beamer が内部でロードするため不要
% subfigure は非推奨（TeX Live が警告を出す）。
\usepackage{subfigure}

% === 参考文献 ================================================
\usepackage{natbib}

% === 欧文フォント =============================================
% newpxtext: Palatino 改良版の本文フォント（mathpazo の現代的後継）
% newpxmath: Palatino 系数式フォント（記号セットが mathpazo より充実）
% \usefonttheme{serif} と組み合わせて Beamer 全体をセリフ体に統一
%
% 代替: Times 系にする場合は以下に切り替え
%   \usepackage{newtxtext}
%   \usepackage{newtxmath}
\usepackage{newpxtext}
\usepackage{newpxmath}

% === 日本語 ==================================================
% otf / pxchfon は \documentclass 直後のメインファイルに記述する想定
% （Beamer のクラスオプション dvipdfmx より先に読まれる必要があるため）
%
%   メインファイルの例:
%     \documentclass[dvipdfmx, t, 10pt]{beamer}
%     \usepackage[deluxe]{otf}
%     \usepackage[hiragino-pron]{pxchfon}
%     \input{my_preamble}

% hyperref の日本語対応（dvipdfmx 環境でのしおり文字化け防止）
\usepackage{pxjahyper}

% === フォントテーマ ==========================================
% 和文: \kanjifamilydefault を \gtdefault にしているのでゴシック体
% 欧文: \usefonttheme{serif} で Palatino（mathpazo）
% → 和文ゴシック + 欧文セリフ の組み合わせ（プレゼン標準）
% 和文も明朝にしたい場合は \gtdefault を \mcdefault に変更する
\renewcommand{\kanjifamilydefault}{\gtdefault}
\usefonttheme{serif}

% 行間（1.0 = デフォルト、1.1〜1.2 が「少し広め」の目安）
\linespread{1.15}

% upLaTeX のゴシック体には b（bold）ウェイトが未定義で、
% \bfseries 使用時に警告が出る。b → bx へのマッピングを明示的に
% 宣言することで警告を抑制する（JY2: 横組み、JT2: 縦組み）
\DeclareFontShape{JY2}{hgt}{b}{n}{<->ssub*hgt/bx/n}{}
\DeclareFontShape{JT2}{hgt}{b}{n}{<->ssub*hgt/bx/n}{}

% === Beamer テーマ ===========================================
\usetheme{default}
\usecolortheme{default}

% ナビゲーションバーを削除
\setbeamertemplate{navigation symbols}{}

% === カラー（モノクロ） =======================================
\setbeamercolor{normal text}{fg=black, bg=white}
\setbeamercolor{headline}{fg=black, bg=white}
\setbeamercolor{footline}{fg=black, bg=white}
\setbeamercolor{frametitle}{fg=black, bg=white}
\setbeamercolor{block title}{fg=black, bg=white}
\setbeamercolor{block body}{fg=black, bg=white}
\setbeamercolor{itemize item}{fg=black}
\setbeamercolor{itemize subitem}{fg=black}
\setbeamercolor{itemize subsubitem}{fg=black}
\setbeamercolor{enumerate item}{fg=black}
\setbeamercolor{enumerate subitem}{fg=black}
\setbeamercolor{enumerate subsubitem}{fg=black}
\setbeamercolor{alerted text}{fg=black}

% ブロックの枠をなくす
\setbeamertemplate{blocks}[default]

% === 余白 ====================================================
\setbeamersize{text margin left=5mm, text margin right=5mm}

% === メタデータ（各ファイルで上書き可） ======================
\newcommand{\coursename}{TikZ：経済学でよく使う図}
\author[山田 知明]{山田 知明}
\title{TikZ：経済学でよく使う図}
\institute[明治大学]{明治大学 商学部}
\date{2026年4月7日}

% === ヘッダー（奇数: 左上に講義名 / 偶数: 右上に講義者名） ==
% leftskip/rightskip でマージンを beamercolorbox 内に閉じ込める
\setbeamertemplate{headline}{%
  \begin{beamercolorbox}[%
    wd=\paperwidth, ht=4mm, dp=1.5mm,
    leftskip=5mm, rightskip=5mm]{headline}%
    \ifodd\insertframenumber
      {\footnotesize\coursename}\hfill
    \else
      \hfill
      {\footnotesize\insertshortauthor%
        \ifx\insertshortinstitute\@empty\else
          \ (\insertshortinstitute)%
        \fi}%
    \fi
  \end{beamercolorbox}%
  % 線：テキスト領域幅（左右 5mm マージン分を差し引く）
  \hbox{\hspace{5mm}\rule{\dimexpr\paperwidth-10mm\relax}{0.1pt}}%
}

% === フッター（右下にページ番号） ============================
% [plain] フレーム（タイトルページ等）では自動的に非表示になる
\setbeamertemplate{footline}{%
  \hbox{\hspace{5mm}\rule{\dimexpr\paperwidth-10mm\relax}{0.1pt}}%
  \begin{beamercolorbox}[%
    wd=\paperwidth, ht=2.5mm, dp=1.5mm,
    leftskip=5mm, rightskip=5mm]{footline}%
    %\hfill{\footnotesize\insertframenumber}%
    \hfill{{\footnotesize\insertframenumber\,/\,\inserttotalframenumber}}
  \end{beamercolorbox}%
}

% === フレームタイトル =========================================
\setbeamertemplate{frametitle}{%
  \vspace{1mm}%
  {\large\bfseries\insertframetitle}%
  \ifx\insertframesubtitle\@empty\else
    \\[0.5mm]{\small\itshape\insertframesubtitle}%
  \fi
  \vspace{-1mm}%
}

% === 箇条書き記号 =============================================
\setbeamertemplate{itemize item}{$\bullet$}
\setbeamertemplate{itemize subitem}{$\circ$}
\setbeamertemplate{itemize subsubitem}{--}

% === タイトルページ ===========================================
\setbeamertemplate{title page}{%
  \vfill
  \centering
  {\LARGE\bfseries\inserttitle}\\[5mm]
  {\normalsize\insertauthor}\\[2mm]
  {\normalsize\insertinstitute}\\[1mm]
  {\normalsize\insertdate}%
  \vfill
}

% === 参考文献スタイル =========================================
\bibliographystyle{ecta}

%
%  前提: \documentclass[dvipdfmx, ...]{beamer}
% ============================================================

% === 数学・記号 ==============================================
% amsmath / amsfonts は Beamer が内部でロード済みなので不要
% amssymb のみ Beamer が読まないため明示的にロードする
\usepackage{amssymb}

% === 図版・作図 ==============================================
% standalone: \input で standalone クラスのファイルを取り込む際に
%   プリアンブル（\documentclass 〜 \begin{document}）を自動的に無視し、
%   本文（TikZ コード等）だけを展開する。
%   → 各 TikZ ファイルを単体コンパイルしながら Beamer にも \input できる。
\usepackage{standalone}

\usepackage{tikz}
\usetikzlibrary{arrows.meta}               % Stealth 矢印等
\usetikzlibrary{decorations.pathreplacing} % 波括弧（DiD図等）
\usetikzlibrary{shapes.geometric, positioning}          % 楕円ノード等（因果推論図）

% PGFPlots: 関数プロット・データ可視化
\usepackage{pgfplots}
\pgfplotsset{compat=1.18}
\usepgfplotslibrary{groupplots}       % 複数パネル（上下・左右配置）
\usepgfplotslibrary{fillbetween}      % 2曲線間の塗りつぶし

% === TikZ 図共通カラー ========================================
% standalone プリアンブルは \input 時に無視されるため、
% 各 TikZ ファイルで使う色はここで定義する。
\definecolor{cfcol}  {RGB}{232, 232, 232}     % 観測不可（薄灰）
\definecolor{cfgray}  {RGB}{100, 140, 200}    % 反事実（薄い青）
\definecolor{expgreen} {RGB}{200, 230, 200}   % 拡張局面シェード
\definecolor{fillcol}{RGB}{180, 215, 255}     % 政策効果の差分領域
\definecolor{gdpblue}  {RGB}{30,  90,  180}   % GDP 線・AD 曲線
\definecolor{goldcol}  {RGB}{190, 140,  10}
\definecolor{recred}   {RGB}{240, 200, 200}   % 後退局面シェード
\definecolor{srasgreen}{RGB}{30,  140, 60}    % SRAS（短期総供給曲線）
\definecolor{trendred} {RGB}{190, 40,  40}    % トレンド線・LRAS
\definecolor{obscol} {RGB}{210, 228, 255}     % 観測可能（薄青）
\definecolor{watercol}{RGB}{175, 210, 240}

% graphicx は Beamer が内部でロードするため不要
% subfigure は非推奨（TeX Live が警告を出す）。
\usepackage{subfigure}

% === 参考文献 ================================================
\usepackage{natbib}

% === 欧文フォント =============================================
% newpxtext: Palatino 改良版の本文フォント（mathpazo の現代的後継）
% newpxmath: Palatino 系数式フォント（記号セットが mathpazo より充実）
% \usefonttheme{serif} と組み合わせて Beamer 全体をセリフ体に統一
%
% 代替: Times 系にする場合は以下に切り替え
%   \usepackage{newtxtext}
%   \usepackage{newtxmath}
\usepackage{newpxtext}
\usepackage{newpxmath}

% === 日本語 ==================================================
% otf / pxchfon は \documentclass 直後のメインファイルに記述する想定
% （Beamer のクラスオプション dvipdfmx より先に読まれる必要があるため）
%
%   メインファイルの例:
%     \documentclass[dvipdfmx, t, 10pt]{beamer}
%     \usepackage[deluxe]{otf}
%     \usepackage[hiragino-pron]{pxchfon}
%     % ============================================================
%  my_preamble.sty
%  upLaTeX + dvipdfmx + Beamer 用プリアンブル
%
%  使い方（\documentclass の直後に）:
%    \input{my_preamble}
%
%  前提: \documentclass[dvipdfmx, ...]{beamer}
% ============================================================

% === 数学・記号 ==============================================
% amsmath / amsfonts は Beamer が内部でロード済みなので不要
% amssymb のみ Beamer が読まないため明示的にロードする
\usepackage{amssymb}

% === 図版・作図 ==============================================
% standalone: \input で standalone クラスのファイルを取り込む際に
%   プリアンブル（\documentclass 〜 \begin{document}）を自動的に無視し、
%   本文（TikZ コード等）だけを展開する。
%   → 各 TikZ ファイルを単体コンパイルしながら Beamer にも \input できる。
\usepackage{standalone}

\usepackage{tikz}
\usetikzlibrary{arrows.meta}               % Stealth 矢印等
\usetikzlibrary{decorations.pathreplacing} % 波括弧（DiD図等）
\usetikzlibrary{shapes.geometric, positioning}          % 楕円ノード等（因果推論図）

% PGFPlots: 関数プロット・データ可視化
\usepackage{pgfplots}
\pgfplotsset{compat=1.18}
\usepgfplotslibrary{groupplots}       % 複数パネル（上下・左右配置）
\usepgfplotslibrary{fillbetween}      % 2曲線間の塗りつぶし

% === TikZ 図共通カラー ========================================
% standalone プリアンブルは \input 時に無視されるため、
% 各 TikZ ファイルで使う色はここで定義する。
\definecolor{cfcol}  {RGB}{232, 232, 232}     % 観測不可（薄灰）
\definecolor{cfgray}  {RGB}{100, 140, 200}    % 反事実（薄い青）
\definecolor{expgreen} {RGB}{200, 230, 200}   % 拡張局面シェード
\definecolor{fillcol}{RGB}{180, 215, 255}     % 政策効果の差分領域
\definecolor{gdpblue}  {RGB}{30,  90,  180}   % GDP 線・AD 曲線
\definecolor{goldcol}  {RGB}{190, 140,  10}
\definecolor{recred}   {RGB}{240, 200, 200}   % 後退局面シェード
\definecolor{srasgreen}{RGB}{30,  140, 60}    % SRAS（短期総供給曲線）
\definecolor{trendred} {RGB}{190, 40,  40}    % トレンド線・LRAS
\definecolor{obscol} {RGB}{210, 228, 255}     % 観測可能（薄青）
\definecolor{watercol}{RGB}{175, 210, 240}

% graphicx は Beamer が内部でロードするため不要
% subfigure は非推奨（TeX Live が警告を出す）。
\usepackage{subfigure}

% === 参考文献 ================================================
\usepackage{natbib}

% === 欧文フォント =============================================
% newpxtext: Palatino 改良版の本文フォント（mathpazo の現代的後継）
% newpxmath: Palatino 系数式フォント（記号セットが mathpazo より充実）
% \usefonttheme{serif} と組み合わせて Beamer 全体をセリフ体に統一
%
% 代替: Times 系にする場合は以下に切り替え
%   \usepackage{newtxtext}
%   \usepackage{newtxmath}
\usepackage{newpxtext}
\usepackage{newpxmath}

% === 日本語 ==================================================
% otf / pxchfon は \documentclass 直後のメインファイルに記述する想定
% （Beamer のクラスオプション dvipdfmx より先に読まれる必要があるため）
%
%   メインファイルの例:
%     \documentclass[dvipdfmx, t, 10pt]{beamer}
%     \usepackage[deluxe]{otf}
%     \usepackage[hiragino-pron]{pxchfon}
%     \input{my_preamble}

% hyperref の日本語対応（dvipdfmx 環境でのしおり文字化け防止）
\usepackage{pxjahyper}

% === フォントテーマ ==========================================
% 和文: \kanjifamilydefault を \gtdefault にしているのでゴシック体
% 欧文: \usefonttheme{serif} で Palatino（mathpazo）
% → 和文ゴシック + 欧文セリフ の組み合わせ（プレゼン標準）
% 和文も明朝にしたい場合は \gtdefault を \mcdefault に変更する
\renewcommand{\kanjifamilydefault}{\gtdefault}
\usefonttheme{serif}

% 行間（1.0 = デフォルト、1.1〜1.2 が「少し広め」の目安）
\linespread{1.15}

% upLaTeX のゴシック体には b（bold）ウェイトが未定義で、
% \bfseries 使用時に警告が出る。b → bx へのマッピングを明示的に
% 宣言することで警告を抑制する（JY2: 横組み、JT2: 縦組み）
\DeclareFontShape{JY2}{hgt}{b}{n}{<->ssub*hgt/bx/n}{}
\DeclareFontShape{JT2}{hgt}{b}{n}{<->ssub*hgt/bx/n}{}

% === Beamer テーマ ===========================================
\usetheme{default}
\usecolortheme{default}

% ナビゲーションバーを削除
\setbeamertemplate{navigation symbols}{}

% === カラー（モノクロ） =======================================
\setbeamercolor{normal text}{fg=black, bg=white}
\setbeamercolor{headline}{fg=black, bg=white}
\setbeamercolor{footline}{fg=black, bg=white}
\setbeamercolor{frametitle}{fg=black, bg=white}
\setbeamercolor{block title}{fg=black, bg=white}
\setbeamercolor{block body}{fg=black, bg=white}
\setbeamercolor{itemize item}{fg=black}
\setbeamercolor{itemize subitem}{fg=black}
\setbeamercolor{itemize subsubitem}{fg=black}
\setbeamercolor{enumerate item}{fg=black}
\setbeamercolor{enumerate subitem}{fg=black}
\setbeamercolor{enumerate subsubitem}{fg=black}
\setbeamercolor{alerted text}{fg=black}

% ブロックの枠をなくす
\setbeamertemplate{blocks}[default]

% === 余白 ====================================================
\setbeamersize{text margin left=5mm, text margin right=5mm}

% === メタデータ（各ファイルで上書き可） ======================
\newcommand{\coursename}{TikZ：経済学でよく使う図}
\author[山田 知明]{山田 知明}
\title{TikZ：経済学でよく使う図}
\institute[明治大学]{明治大学 商学部}
\date{2026年4月7日}

% === ヘッダー（奇数: 左上に講義名 / 偶数: 右上に講義者名） ==
% leftskip/rightskip でマージンを beamercolorbox 内に閉じ込める
\setbeamertemplate{headline}{%
  \begin{beamercolorbox}[%
    wd=\paperwidth, ht=4mm, dp=1.5mm,
    leftskip=5mm, rightskip=5mm]{headline}%
    \ifodd\insertframenumber
      {\footnotesize\coursename}\hfill
    \else
      \hfill
      {\footnotesize\insertshortauthor%
        \ifx\insertshortinstitute\@empty\else
          \ (\insertshortinstitute)%
        \fi}%
    \fi
  \end{beamercolorbox}%
  % 線：テキスト領域幅（左右 5mm マージン分を差し引く）
  \hbox{\hspace{5mm}\rule{\dimexpr\paperwidth-10mm\relax}{0.1pt}}%
}

% === フッター（右下にページ番号） ============================
% [plain] フレーム（タイトルページ等）では自動的に非表示になる
\setbeamertemplate{footline}{%
  \hbox{\hspace{5mm}\rule{\dimexpr\paperwidth-10mm\relax}{0.1pt}}%
  \begin{beamercolorbox}[%
    wd=\paperwidth, ht=2.5mm, dp=1.5mm,
    leftskip=5mm, rightskip=5mm]{footline}%
    %\hfill{\footnotesize\insertframenumber}%
    \hfill{{\footnotesize\insertframenumber\,/\,\inserttotalframenumber}}
  \end{beamercolorbox}%
}

% === フレームタイトル =========================================
\setbeamertemplate{frametitle}{%
  \vspace{1mm}%
  {\large\bfseries\insertframetitle}%
  \ifx\insertframesubtitle\@empty\else
    \\[0.5mm]{\small\itshape\insertframesubtitle}%
  \fi
  \vspace{-1mm}%
}

% === 箇条書き記号 =============================================
\setbeamertemplate{itemize item}{$\bullet$}
\setbeamertemplate{itemize subitem}{$\circ$}
\setbeamertemplate{itemize subsubitem}{--}

% === タイトルページ ===========================================
\setbeamertemplate{title page}{%
  \vfill
  \centering
  {\LARGE\bfseries\inserttitle}\\[5mm]
  {\normalsize\insertauthor}\\[2mm]
  {\normalsize\insertinstitute}\\[1mm]
  {\normalsize\insertdate}%
  \vfill
}

% === 参考文献スタイル =========================================
\bibliographystyle{ecta}


% hyperref の日本語対応（dvipdfmx 環境でのしおり文字化け防止）
\usepackage{pxjahyper}

% === フォントテーマ ==========================================
% 和文: \kanjifamilydefault を \gtdefault にしているのでゴシック体
% 欧文: \usefonttheme{serif} で Palatino（mathpazo）
% → 和文ゴシック + 欧文セリフ の組み合わせ（プレゼン標準）
% 和文も明朝にしたい場合は \gtdefault を \mcdefault に変更する
\renewcommand{\kanjifamilydefault}{\gtdefault}
\usefonttheme{serif}

% 行間（1.0 = デフォルト、1.1〜1.2 が「少し広め」の目安）
\linespread{1.15}

% upLaTeX のゴシック体には b（bold）ウェイトが未定義で、
% \bfseries 使用時に警告が出る。b → bx へのマッピングを明示的に
% 宣言することで警告を抑制する（JY2: 横組み、JT2: 縦組み）
\DeclareFontShape{JY2}{hgt}{b}{n}{<->ssub*hgt/bx/n}{}
\DeclareFontShape{JT2}{hgt}{b}{n}{<->ssub*hgt/bx/n}{}

% === Beamer テーマ ===========================================
\usetheme{default}
\usecolortheme{default}

% ナビゲーションバーを削除
\setbeamertemplate{navigation symbols}{}

% === カラー（モノクロ） =======================================
\setbeamercolor{normal text}{fg=black, bg=white}
\setbeamercolor{headline}{fg=black, bg=white}
\setbeamercolor{footline}{fg=black, bg=white}
\setbeamercolor{frametitle}{fg=black, bg=white}
\setbeamercolor{block title}{fg=black, bg=white}
\setbeamercolor{block body}{fg=black, bg=white}
\setbeamercolor{itemize item}{fg=black}
\setbeamercolor{itemize subitem}{fg=black}
\setbeamercolor{itemize subsubitem}{fg=black}
\setbeamercolor{enumerate item}{fg=black}
\setbeamercolor{enumerate subitem}{fg=black}
\setbeamercolor{enumerate subsubitem}{fg=black}
\setbeamercolor{alerted text}{fg=black}

% ブロックの枠をなくす
\setbeamertemplate{blocks}[default]

% === 余白 ====================================================
\setbeamersize{text margin left=5mm, text margin right=5mm}

% === メタデータ（各ファイルで上書き可） ======================
\newcommand{\coursename}{TikZ：経済学でよく使う図}
\author[山田 知明]{山田 知明}
\title{TikZ：経済学でよく使う図}
\institute[明治大学]{明治大学 商学部}
\date{2026年4月7日}

% === ヘッダー（奇数: 左上に講義名 / 偶数: 右上に講義者名） ==
% leftskip/rightskip でマージンを beamercolorbox 内に閉じ込める
\setbeamertemplate{headline}{%
  \begin{beamercolorbox}[%
    wd=\paperwidth, ht=4mm, dp=1.5mm,
    leftskip=5mm, rightskip=5mm]{headline}%
    \ifodd\insertframenumber
      {\footnotesize\coursename}\hfill
    \else
      \hfill
      {\footnotesize\insertshortauthor%
        \ifx\insertshortinstitute\@empty\else
          \ (\insertshortinstitute)%
        \fi}%
    \fi
  \end{beamercolorbox}%
  % 線：テキスト領域幅（左右 5mm マージン分を差し引く）
  \hbox{\hspace{5mm}\rule{\dimexpr\paperwidth-10mm\relax}{0.1pt}}%
}

% === フッター（右下にページ番号） ============================
% [plain] フレーム（タイトルページ等）では自動的に非表示になる
\setbeamertemplate{footline}{%
  \hbox{\hspace{5mm}\rule{\dimexpr\paperwidth-10mm\relax}{0.1pt}}%
  \begin{beamercolorbox}[%
    wd=\paperwidth, ht=2.5mm, dp=1.5mm,
    leftskip=5mm, rightskip=5mm]{footline}%
    %\hfill{\footnotesize\insertframenumber}%
    \hfill{{\footnotesize\insertframenumber\,/\,\inserttotalframenumber}}
  \end{beamercolorbox}%
}

% === フレームタイトル =========================================
\setbeamertemplate{frametitle}{%
  \vspace{1mm}%
  {\large\bfseries\insertframetitle}%
  \ifx\insertframesubtitle\@empty\else
    \\[0.5mm]{\small\itshape\insertframesubtitle}%
  \fi
  \vspace{-1mm}%
}

% === 箇条書き記号 =============================================
\setbeamertemplate{itemize item}{$\bullet$}
\setbeamertemplate{itemize subitem}{$\circ$}
\setbeamertemplate{itemize subsubitem}{--}

% === タイトルページ ===========================================
\setbeamertemplate{title page}{%
  \vfill
  \centering
  {\LARGE\bfseries\inserttitle}\\[5mm]
  {\normalsize\insertauthor}\\[2mm]
  {\normalsize\insertinstitute}\\[1mm]
  {\normalsize\insertdate}%
  \vfill
}

% === 参考文献スタイル =========================================
\bibliographystyle{ecta}

%
%  前提: \documentclass[dvipdfmx, ...]{beamer}
% ============================================================

% === 数学・記号 ==============================================
% amsmath / amsfonts は Beamer が内部でロード済みなので不要
% amssymb のみ Beamer が読まないため明示的にロードする
\usepackage{amssymb}

% === 図版・作図 ==============================================
% standalone: \input で standalone クラスのファイルを取り込む際に
%   プリアンブル（\documentclass 〜 \begin{document}）を自動的に無視し、
%   本文（TikZ コード等）だけを展開する。
%   → 各 TikZ ファイルを単体コンパイルしながら Beamer にも \input できる。
\usepackage{standalone}

\usepackage{tikz}
\usetikzlibrary{arrows.meta}               % Stealth 矢印等
\usetikzlibrary{decorations.pathreplacing} % 波括弧（DiD図等）
\usetikzlibrary{shapes.geometric, positioning}          % 楕円ノード等（因果推論図）

% PGFPlots: 関数プロット・データ可視化
\usepackage{pgfplots}
\pgfplotsset{compat=1.18}
\usepgfplotslibrary{groupplots}       % 複数パネル（上下・左右配置）
\usepgfplotslibrary{fillbetween}      % 2曲線間の塗りつぶし

% === TikZ 図共通カラー ========================================
% standalone プリアンブルは \input 時に無視されるため、
% 各 TikZ ファイルで使う色はここで定義する。
\definecolor{cfcol}  {RGB}{232, 232, 232}     % 観測不可（薄灰）
\definecolor{cfgray}  {RGB}{100, 140, 200}    % 反事実（薄い青）
\definecolor{expgreen} {RGB}{200, 230, 200}   % 拡張局面シェード
\definecolor{fillcol}{RGB}{180, 215, 255}     % 政策効果の差分領域
\definecolor{gdpblue}  {RGB}{30,  90,  180}   % GDP 線・AD 曲線
\definecolor{goldcol}  {RGB}{190, 140,  10}
\definecolor{recred}   {RGB}{240, 200, 200}   % 後退局面シェード
\definecolor{srasgreen}{RGB}{30,  140, 60}    % SRAS（短期総供給曲線）
\definecolor{trendred} {RGB}{190, 40,  40}    % トレンド線・LRAS
\definecolor{obscol} {RGB}{210, 228, 255}     % 観測可能（薄青）
\definecolor{watercol}{RGB}{175, 210, 240}

% graphicx は Beamer が内部でロードするため不要
% subfigure は非推奨（TeX Live が警告を出す）。
\usepackage{subfigure}

% === 参考文献 ================================================
\usepackage{natbib}

% === 欧文フォント =============================================
% newpxtext: Palatino 改良版の本文フォント（mathpazo の現代的後継）
% newpxmath: Palatino 系数式フォント（記号セットが mathpazo より充実）
% \usefonttheme{serif} と組み合わせて Beamer 全体をセリフ体に統一
%
% 代替: Times 系にする場合は以下に切り替え
%   \usepackage{newtxtext}
%   \usepackage{newtxmath}
\usepackage{newpxtext}
\usepackage{newpxmath}

% === 日本語 ==================================================
% otf / pxchfon は \documentclass 直後のメインファイルに記述する想定
% （Beamer のクラスオプション dvipdfmx より先に読まれる必要があるため）
%
%   メインファイルの例:
%     \documentclass[dvipdfmx, t, 10pt]{beamer}
%     \usepackage[deluxe]{otf}
%     \usepackage[hiragino-pron]{pxchfon}
%     % ============================================================
%  my_preamble.sty
%  upLaTeX + dvipdfmx + Beamer 用プリアンブル
%
%  使い方（\documentclass の直後に）:
%    % ============================================================
%  my_preamble.sty
%  upLaTeX + dvipdfmx + Beamer 用プリアンブル
%
%  使い方（\documentclass の直後に）:
%    \input{my_preamble}
%
%  前提: \documentclass[dvipdfmx, ...]{beamer}
% ============================================================

% === 数学・記号 ==============================================
% amsmath / amsfonts は Beamer が内部でロード済みなので不要
% amssymb のみ Beamer が読まないため明示的にロードする
\usepackage{amssymb}

% === 図版・作図 ==============================================
% standalone: \input で standalone クラスのファイルを取り込む際に
%   プリアンブル（\documentclass 〜 \begin{document}）を自動的に無視し、
%   本文（TikZ コード等）だけを展開する。
%   → 各 TikZ ファイルを単体コンパイルしながら Beamer にも \input できる。
\usepackage{standalone}

\usepackage{tikz}
\usetikzlibrary{arrows.meta}               % Stealth 矢印等
\usetikzlibrary{decorations.pathreplacing} % 波括弧（DiD図等）
\usetikzlibrary{shapes.geometric, positioning}          % 楕円ノード等（因果推論図）

% PGFPlots: 関数プロット・データ可視化
\usepackage{pgfplots}
\pgfplotsset{compat=1.18}
\usepgfplotslibrary{groupplots}       % 複数パネル（上下・左右配置）
\usepgfplotslibrary{fillbetween}      % 2曲線間の塗りつぶし

% === TikZ 図共通カラー ========================================
% standalone プリアンブルは \input 時に無視されるため、
% 各 TikZ ファイルで使う色はここで定義する。
\definecolor{cfcol}  {RGB}{232, 232, 232}     % 観測不可（薄灰）
\definecolor{cfgray}  {RGB}{100, 140, 200}    % 反事実（薄い青）
\definecolor{expgreen} {RGB}{200, 230, 200}   % 拡張局面シェード
\definecolor{fillcol}{RGB}{180, 215, 255}     % 政策効果の差分領域
\definecolor{gdpblue}  {RGB}{30,  90,  180}   % GDP 線・AD 曲線
\definecolor{goldcol}  {RGB}{190, 140,  10}
\definecolor{recred}   {RGB}{240, 200, 200}   % 後退局面シェード
\definecolor{srasgreen}{RGB}{30,  140, 60}    % SRAS（短期総供給曲線）
\definecolor{trendred} {RGB}{190, 40,  40}    % トレンド線・LRAS
\definecolor{obscol} {RGB}{210, 228, 255}     % 観測可能（薄青）
\definecolor{watercol}{RGB}{175, 210, 240}

% graphicx は Beamer が内部でロードするため不要
% subfigure は非推奨（TeX Live が警告を出す）。
\usepackage{subfigure}

% === 参考文献 ================================================
\usepackage{natbib}

% === 欧文フォント =============================================
% newpxtext: Palatino 改良版の本文フォント（mathpazo の現代的後継）
% newpxmath: Palatino 系数式フォント（記号セットが mathpazo より充実）
% \usefonttheme{serif} と組み合わせて Beamer 全体をセリフ体に統一
%
% 代替: Times 系にする場合は以下に切り替え
%   \usepackage{newtxtext}
%   \usepackage{newtxmath}
\usepackage{newpxtext}
\usepackage{newpxmath}

% === 日本語 ==================================================
% otf / pxchfon は \documentclass 直後のメインファイルに記述する想定
% （Beamer のクラスオプション dvipdfmx より先に読まれる必要があるため）
%
%   メインファイルの例:
%     \documentclass[dvipdfmx, t, 10pt]{beamer}
%     \usepackage[deluxe]{otf}
%     \usepackage[hiragino-pron]{pxchfon}
%     \input{my_preamble}

% hyperref の日本語対応（dvipdfmx 環境でのしおり文字化け防止）
\usepackage{pxjahyper}

% === フォントテーマ ==========================================
% 和文: \kanjifamilydefault を \gtdefault にしているのでゴシック体
% 欧文: \usefonttheme{serif} で Palatino（mathpazo）
% → 和文ゴシック + 欧文セリフ の組み合わせ（プレゼン標準）
% 和文も明朝にしたい場合は \gtdefault を \mcdefault に変更する
\renewcommand{\kanjifamilydefault}{\gtdefault}
\usefonttheme{serif}

% 行間（1.0 = デフォルト、1.1〜1.2 が「少し広め」の目安）
\linespread{1.15}

% upLaTeX のゴシック体には b（bold）ウェイトが未定義で、
% \bfseries 使用時に警告が出る。b → bx へのマッピングを明示的に
% 宣言することで警告を抑制する（JY2: 横組み、JT2: 縦組み）
\DeclareFontShape{JY2}{hgt}{b}{n}{<->ssub*hgt/bx/n}{}
\DeclareFontShape{JT2}{hgt}{b}{n}{<->ssub*hgt/bx/n}{}

% === Beamer テーマ ===========================================
\usetheme{default}
\usecolortheme{default}

% ナビゲーションバーを削除
\setbeamertemplate{navigation symbols}{}

% === カラー（モノクロ） =======================================
\setbeamercolor{normal text}{fg=black, bg=white}
\setbeamercolor{headline}{fg=black, bg=white}
\setbeamercolor{footline}{fg=black, bg=white}
\setbeamercolor{frametitle}{fg=black, bg=white}
\setbeamercolor{block title}{fg=black, bg=white}
\setbeamercolor{block body}{fg=black, bg=white}
\setbeamercolor{itemize item}{fg=black}
\setbeamercolor{itemize subitem}{fg=black}
\setbeamercolor{itemize subsubitem}{fg=black}
\setbeamercolor{enumerate item}{fg=black}
\setbeamercolor{enumerate subitem}{fg=black}
\setbeamercolor{enumerate subsubitem}{fg=black}
\setbeamercolor{alerted text}{fg=black}

% ブロックの枠をなくす
\setbeamertemplate{blocks}[default]

% === 余白 ====================================================
\setbeamersize{text margin left=5mm, text margin right=5mm}

% === メタデータ（各ファイルで上書き可） ======================
\newcommand{\coursename}{TikZ：経済学でよく使う図}
\author[山田 知明]{山田 知明}
\title{TikZ：経済学でよく使う図}
\institute[明治大学]{明治大学 商学部}
\date{2026年4月7日}

% === ヘッダー（奇数: 左上に講義名 / 偶数: 右上に講義者名） ==
% leftskip/rightskip でマージンを beamercolorbox 内に閉じ込める
\setbeamertemplate{headline}{%
  \begin{beamercolorbox}[%
    wd=\paperwidth, ht=4mm, dp=1.5mm,
    leftskip=5mm, rightskip=5mm]{headline}%
    \ifodd\insertframenumber
      {\footnotesize\coursename}\hfill
    \else
      \hfill
      {\footnotesize\insertshortauthor%
        \ifx\insertshortinstitute\@empty\else
          \ (\insertshortinstitute)%
        \fi}%
    \fi
  \end{beamercolorbox}%
  % 線：テキスト領域幅（左右 5mm マージン分を差し引く）
  \hbox{\hspace{5mm}\rule{\dimexpr\paperwidth-10mm\relax}{0.1pt}}%
}

% === フッター（右下にページ番号） ============================
% [plain] フレーム（タイトルページ等）では自動的に非表示になる
\setbeamertemplate{footline}{%
  \hbox{\hspace{5mm}\rule{\dimexpr\paperwidth-10mm\relax}{0.1pt}}%
  \begin{beamercolorbox}[%
    wd=\paperwidth, ht=2.5mm, dp=1.5mm,
    leftskip=5mm, rightskip=5mm]{footline}%
    %\hfill{\footnotesize\insertframenumber}%
    \hfill{{\footnotesize\insertframenumber\,/\,\inserttotalframenumber}}
  \end{beamercolorbox}%
}

% === フレームタイトル =========================================
\setbeamertemplate{frametitle}{%
  \vspace{1mm}%
  {\large\bfseries\insertframetitle}%
  \ifx\insertframesubtitle\@empty\else
    \\[0.5mm]{\small\itshape\insertframesubtitle}%
  \fi
  \vspace{-1mm}%
}

% === 箇条書き記号 =============================================
\setbeamertemplate{itemize item}{$\bullet$}
\setbeamertemplate{itemize subitem}{$\circ$}
\setbeamertemplate{itemize subsubitem}{--}

% === タイトルページ ===========================================
\setbeamertemplate{title page}{%
  \vfill
  \centering
  {\LARGE\bfseries\inserttitle}\\[5mm]
  {\normalsize\insertauthor}\\[2mm]
  {\normalsize\insertinstitute}\\[1mm]
  {\normalsize\insertdate}%
  \vfill
}

% === 参考文献スタイル =========================================
\bibliographystyle{ecta}

%
%  前提: \documentclass[dvipdfmx, ...]{beamer}
% ============================================================

% === 数学・記号 ==============================================
% amsmath / amsfonts は Beamer が内部でロード済みなので不要
% amssymb のみ Beamer が読まないため明示的にロードする
\usepackage{amssymb}

% === 図版・作図 ==============================================
% standalone: \input で standalone クラスのファイルを取り込む際に
%   プリアンブル（\documentclass 〜 \begin{document}）を自動的に無視し、
%   本文（TikZ コード等）だけを展開する。
%   → 各 TikZ ファイルを単体コンパイルしながら Beamer にも \input できる。
\usepackage{standalone}

\usepackage{tikz}
\usetikzlibrary{arrows.meta}               % Stealth 矢印等
\usetikzlibrary{decorations.pathreplacing} % 波括弧（DiD図等）
\usetikzlibrary{shapes.geometric, positioning}          % 楕円ノード等（因果推論図）

% PGFPlots: 関数プロット・データ可視化
\usepackage{pgfplots}
\pgfplotsset{compat=1.18}
\usepgfplotslibrary{groupplots}       % 複数パネル（上下・左右配置）
\usepgfplotslibrary{fillbetween}      % 2曲線間の塗りつぶし

% === TikZ 図共通カラー ========================================
% standalone プリアンブルは \input 時に無視されるため、
% 各 TikZ ファイルで使う色はここで定義する。
\definecolor{cfcol}  {RGB}{232, 232, 232}     % 観測不可（薄灰）
\definecolor{cfgray}  {RGB}{100, 140, 200}    % 反事実（薄い青）
\definecolor{expgreen} {RGB}{200, 230, 200}   % 拡張局面シェード
\definecolor{fillcol}{RGB}{180, 215, 255}     % 政策効果の差分領域
\definecolor{gdpblue}  {RGB}{30,  90,  180}   % GDP 線・AD 曲線
\definecolor{goldcol}  {RGB}{190, 140,  10}
\definecolor{recred}   {RGB}{240, 200, 200}   % 後退局面シェード
\definecolor{srasgreen}{RGB}{30,  140, 60}    % SRAS（短期総供給曲線）
\definecolor{trendred} {RGB}{190, 40,  40}    % トレンド線・LRAS
\definecolor{obscol} {RGB}{210, 228, 255}     % 観測可能（薄青）
\definecolor{watercol}{RGB}{175, 210, 240}

% graphicx は Beamer が内部でロードするため不要
% subfigure は非推奨（TeX Live が警告を出す）。
\usepackage{subfigure}

% === 参考文献 ================================================
\usepackage{natbib}

% === 欧文フォント =============================================
% newpxtext: Palatino 改良版の本文フォント（mathpazo の現代的後継）
% newpxmath: Palatino 系数式フォント（記号セットが mathpazo より充実）
% \usefonttheme{serif} と組み合わせて Beamer 全体をセリフ体に統一
%
% 代替: Times 系にする場合は以下に切り替え
%   \usepackage{newtxtext}
%   \usepackage{newtxmath}
\usepackage{newpxtext}
\usepackage{newpxmath}

% === 日本語 ==================================================
% otf / pxchfon は \documentclass 直後のメインファイルに記述する想定
% （Beamer のクラスオプション dvipdfmx より先に読まれる必要があるため）
%
%   メインファイルの例:
%     \documentclass[dvipdfmx, t, 10pt]{beamer}
%     \usepackage[deluxe]{otf}
%     \usepackage[hiragino-pron]{pxchfon}
%     % ============================================================
%  my_preamble.sty
%  upLaTeX + dvipdfmx + Beamer 用プリアンブル
%
%  使い方（\documentclass の直後に）:
%    \input{my_preamble}
%
%  前提: \documentclass[dvipdfmx, ...]{beamer}
% ============================================================

% === 数学・記号 ==============================================
% amsmath / amsfonts は Beamer が内部でロード済みなので不要
% amssymb のみ Beamer が読まないため明示的にロードする
\usepackage{amssymb}

% === 図版・作図 ==============================================
% standalone: \input で standalone クラスのファイルを取り込む際に
%   プリアンブル（\documentclass 〜 \begin{document}）を自動的に無視し、
%   本文（TikZ コード等）だけを展開する。
%   → 各 TikZ ファイルを単体コンパイルしながら Beamer にも \input できる。
\usepackage{standalone}

\usepackage{tikz}
\usetikzlibrary{arrows.meta}               % Stealth 矢印等
\usetikzlibrary{decorations.pathreplacing} % 波括弧（DiD図等）
\usetikzlibrary{shapes.geometric, positioning}          % 楕円ノード等（因果推論図）

% PGFPlots: 関数プロット・データ可視化
\usepackage{pgfplots}
\pgfplotsset{compat=1.18}
\usepgfplotslibrary{groupplots}       % 複数パネル（上下・左右配置）
\usepgfplotslibrary{fillbetween}      % 2曲線間の塗りつぶし

% === TikZ 図共通カラー ========================================
% standalone プリアンブルは \input 時に無視されるため、
% 各 TikZ ファイルで使う色はここで定義する。
\definecolor{cfcol}  {RGB}{232, 232, 232}     % 観測不可（薄灰）
\definecolor{cfgray}  {RGB}{100, 140, 200}    % 反事実（薄い青）
\definecolor{expgreen} {RGB}{200, 230, 200}   % 拡張局面シェード
\definecolor{fillcol}{RGB}{180, 215, 255}     % 政策効果の差分領域
\definecolor{gdpblue}  {RGB}{30,  90,  180}   % GDP 線・AD 曲線
\definecolor{goldcol}  {RGB}{190, 140,  10}
\definecolor{recred}   {RGB}{240, 200, 200}   % 後退局面シェード
\definecolor{srasgreen}{RGB}{30,  140, 60}    % SRAS（短期総供給曲線）
\definecolor{trendred} {RGB}{190, 40,  40}    % トレンド線・LRAS
\definecolor{obscol} {RGB}{210, 228, 255}     % 観測可能（薄青）
\definecolor{watercol}{RGB}{175, 210, 240}

% graphicx は Beamer が内部でロードするため不要
% subfigure は非推奨（TeX Live が警告を出す）。
\usepackage{subfigure}

% === 参考文献 ================================================
\usepackage{natbib}

% === 欧文フォント =============================================
% newpxtext: Palatino 改良版の本文フォント（mathpazo の現代的後継）
% newpxmath: Palatino 系数式フォント（記号セットが mathpazo より充実）
% \usefonttheme{serif} と組み合わせて Beamer 全体をセリフ体に統一
%
% 代替: Times 系にする場合は以下に切り替え
%   \usepackage{newtxtext}
%   \usepackage{newtxmath}
\usepackage{newpxtext}
\usepackage{newpxmath}

% === 日本語 ==================================================
% otf / pxchfon は \documentclass 直後のメインファイルに記述する想定
% （Beamer のクラスオプション dvipdfmx より先に読まれる必要があるため）
%
%   メインファイルの例:
%     \documentclass[dvipdfmx, t, 10pt]{beamer}
%     \usepackage[deluxe]{otf}
%     \usepackage[hiragino-pron]{pxchfon}
%     \input{my_preamble}

% hyperref の日本語対応（dvipdfmx 環境でのしおり文字化け防止）
\usepackage{pxjahyper}

% === フォントテーマ ==========================================
% 和文: \kanjifamilydefault を \gtdefault にしているのでゴシック体
% 欧文: \usefonttheme{serif} で Palatino（mathpazo）
% → 和文ゴシック + 欧文セリフ の組み合わせ（プレゼン標準）
% 和文も明朝にしたい場合は \gtdefault を \mcdefault に変更する
\renewcommand{\kanjifamilydefault}{\gtdefault}
\usefonttheme{serif}

% 行間（1.0 = デフォルト、1.1〜1.2 が「少し広め」の目安）
\linespread{1.15}

% upLaTeX のゴシック体には b（bold）ウェイトが未定義で、
% \bfseries 使用時に警告が出る。b → bx へのマッピングを明示的に
% 宣言することで警告を抑制する（JY2: 横組み、JT2: 縦組み）
\DeclareFontShape{JY2}{hgt}{b}{n}{<->ssub*hgt/bx/n}{}
\DeclareFontShape{JT2}{hgt}{b}{n}{<->ssub*hgt/bx/n}{}

% === Beamer テーマ ===========================================
\usetheme{default}
\usecolortheme{default}

% ナビゲーションバーを削除
\setbeamertemplate{navigation symbols}{}

% === カラー（モノクロ） =======================================
\setbeamercolor{normal text}{fg=black, bg=white}
\setbeamercolor{headline}{fg=black, bg=white}
\setbeamercolor{footline}{fg=black, bg=white}
\setbeamercolor{frametitle}{fg=black, bg=white}
\setbeamercolor{block title}{fg=black, bg=white}
\setbeamercolor{block body}{fg=black, bg=white}
\setbeamercolor{itemize item}{fg=black}
\setbeamercolor{itemize subitem}{fg=black}
\setbeamercolor{itemize subsubitem}{fg=black}
\setbeamercolor{enumerate item}{fg=black}
\setbeamercolor{enumerate subitem}{fg=black}
\setbeamercolor{enumerate subsubitem}{fg=black}
\setbeamercolor{alerted text}{fg=black}

% ブロックの枠をなくす
\setbeamertemplate{blocks}[default]

% === 余白 ====================================================
\setbeamersize{text margin left=5mm, text margin right=5mm}

% === メタデータ（各ファイルで上書き可） ======================
\newcommand{\coursename}{TikZ：経済学でよく使う図}
\author[山田 知明]{山田 知明}
\title{TikZ：経済学でよく使う図}
\institute[明治大学]{明治大学 商学部}
\date{2026年4月7日}

% === ヘッダー（奇数: 左上に講義名 / 偶数: 右上に講義者名） ==
% leftskip/rightskip でマージンを beamercolorbox 内に閉じ込める
\setbeamertemplate{headline}{%
  \begin{beamercolorbox}[%
    wd=\paperwidth, ht=4mm, dp=1.5mm,
    leftskip=5mm, rightskip=5mm]{headline}%
    \ifodd\insertframenumber
      {\footnotesize\coursename}\hfill
    \else
      \hfill
      {\footnotesize\insertshortauthor%
        \ifx\insertshortinstitute\@empty\else
          \ (\insertshortinstitute)%
        \fi}%
    \fi
  \end{beamercolorbox}%
  % 線：テキスト領域幅（左右 5mm マージン分を差し引く）
  \hbox{\hspace{5mm}\rule{\dimexpr\paperwidth-10mm\relax}{0.1pt}}%
}

% === フッター（右下にページ番号） ============================
% [plain] フレーム（タイトルページ等）では自動的に非表示になる
\setbeamertemplate{footline}{%
  \hbox{\hspace{5mm}\rule{\dimexpr\paperwidth-10mm\relax}{0.1pt}}%
  \begin{beamercolorbox}[%
    wd=\paperwidth, ht=2.5mm, dp=1.5mm,
    leftskip=5mm, rightskip=5mm]{footline}%
    %\hfill{\footnotesize\insertframenumber}%
    \hfill{{\footnotesize\insertframenumber\,/\,\inserttotalframenumber}}
  \end{beamercolorbox}%
}

% === フレームタイトル =========================================
\setbeamertemplate{frametitle}{%
  \vspace{1mm}%
  {\large\bfseries\insertframetitle}%
  \ifx\insertframesubtitle\@empty\else
    \\[0.5mm]{\small\itshape\insertframesubtitle}%
  \fi
  \vspace{-1mm}%
}

% === 箇条書き記号 =============================================
\setbeamertemplate{itemize item}{$\bullet$}
\setbeamertemplate{itemize subitem}{$\circ$}
\setbeamertemplate{itemize subsubitem}{--}

% === タイトルページ ===========================================
\setbeamertemplate{title page}{%
  \vfill
  \centering
  {\LARGE\bfseries\inserttitle}\\[5mm]
  {\normalsize\insertauthor}\\[2mm]
  {\normalsize\insertinstitute}\\[1mm]
  {\normalsize\insertdate}%
  \vfill
}

% === 参考文献スタイル =========================================
\bibliographystyle{ecta}


% hyperref の日本語対応（dvipdfmx 環境でのしおり文字化け防止）
\usepackage{pxjahyper}

% === フォントテーマ ==========================================
% 和文: \kanjifamilydefault を \gtdefault にしているのでゴシック体
% 欧文: \usefonttheme{serif} で Palatino（mathpazo）
% → 和文ゴシック + 欧文セリフ の組み合わせ（プレゼン標準）
% 和文も明朝にしたい場合は \gtdefault を \mcdefault に変更する
\renewcommand{\kanjifamilydefault}{\gtdefault}
\usefonttheme{serif}

% 行間（1.0 = デフォルト、1.1〜1.2 が「少し広め」の目安）
\linespread{1.15}

% upLaTeX のゴシック体には b（bold）ウェイトが未定義で、
% \bfseries 使用時に警告が出る。b → bx へのマッピングを明示的に
% 宣言することで警告を抑制する（JY2: 横組み、JT2: 縦組み）
\DeclareFontShape{JY2}{hgt}{b}{n}{<->ssub*hgt/bx/n}{}
\DeclareFontShape{JT2}{hgt}{b}{n}{<->ssub*hgt/bx/n}{}

% === Beamer テーマ ===========================================
\usetheme{default}
\usecolortheme{default}

% ナビゲーションバーを削除
\setbeamertemplate{navigation symbols}{}

% === カラー（モノクロ） =======================================
\setbeamercolor{normal text}{fg=black, bg=white}
\setbeamercolor{headline}{fg=black, bg=white}
\setbeamercolor{footline}{fg=black, bg=white}
\setbeamercolor{frametitle}{fg=black, bg=white}
\setbeamercolor{block title}{fg=black, bg=white}
\setbeamercolor{block body}{fg=black, bg=white}
\setbeamercolor{itemize item}{fg=black}
\setbeamercolor{itemize subitem}{fg=black}
\setbeamercolor{itemize subsubitem}{fg=black}
\setbeamercolor{enumerate item}{fg=black}
\setbeamercolor{enumerate subitem}{fg=black}
\setbeamercolor{enumerate subsubitem}{fg=black}
\setbeamercolor{alerted text}{fg=black}

% ブロックの枠をなくす
\setbeamertemplate{blocks}[default]

% === 余白 ====================================================
\setbeamersize{text margin left=5mm, text margin right=5mm}

% === メタデータ（各ファイルで上書き可） ======================
\newcommand{\coursename}{TikZ：経済学でよく使う図}
\author[山田 知明]{山田 知明}
\title{TikZ：経済学でよく使う図}
\institute[明治大学]{明治大学 商学部}
\date{2026年4月7日}

% === ヘッダー（奇数: 左上に講義名 / 偶数: 右上に講義者名） ==
% leftskip/rightskip でマージンを beamercolorbox 内に閉じ込める
\setbeamertemplate{headline}{%
  \begin{beamercolorbox}[%
    wd=\paperwidth, ht=4mm, dp=1.5mm,
    leftskip=5mm, rightskip=5mm]{headline}%
    \ifodd\insertframenumber
      {\footnotesize\coursename}\hfill
    \else
      \hfill
      {\footnotesize\insertshortauthor%
        \ifx\insertshortinstitute\@empty\else
          \ (\insertshortinstitute)%
        \fi}%
    \fi
  \end{beamercolorbox}%
  % 線：テキスト領域幅（左右 5mm マージン分を差し引く）
  \hbox{\hspace{5mm}\rule{\dimexpr\paperwidth-10mm\relax}{0.1pt}}%
}

% === フッター（右下にページ番号） ============================
% [plain] フレーム（タイトルページ等）では自動的に非表示になる
\setbeamertemplate{footline}{%
  \hbox{\hspace{5mm}\rule{\dimexpr\paperwidth-10mm\relax}{0.1pt}}%
  \begin{beamercolorbox}[%
    wd=\paperwidth, ht=2.5mm, dp=1.5mm,
    leftskip=5mm, rightskip=5mm]{footline}%
    %\hfill{\footnotesize\insertframenumber}%
    \hfill{{\footnotesize\insertframenumber\,/\,\inserttotalframenumber}}
  \end{beamercolorbox}%
}

% === フレームタイトル =========================================
\setbeamertemplate{frametitle}{%
  \vspace{1mm}%
  {\large\bfseries\insertframetitle}%
  \ifx\insertframesubtitle\@empty\else
    \\[0.5mm]{\small\itshape\insertframesubtitle}%
  \fi
  \vspace{-1mm}%
}

% === 箇条書き記号 =============================================
\setbeamertemplate{itemize item}{$\bullet$}
\setbeamertemplate{itemize subitem}{$\circ$}
\setbeamertemplate{itemize subsubitem}{--}

% === タイトルページ ===========================================
\setbeamertemplate{title page}{%
  \vfill
  \centering
  {\LARGE\bfseries\inserttitle}\\[5mm]
  {\normalsize\insertauthor}\\[2mm]
  {\normalsize\insertinstitute}\\[1mm]
  {\normalsize\insertdate}%
  \vfill
}

% === 参考文献スタイル =========================================
\bibliographystyle{ecta}


% hyperref の日本語対応（dvipdfmx 環境でのしおり文字化け防止）
\usepackage{pxjahyper}

% === フォントテーマ ==========================================
% 和文: \kanjifamilydefault を \gtdefault にしているのでゴシック体
% 欧文: \usefonttheme{serif} で Palatino（mathpazo）
% → 和文ゴシック + 欧文セリフ の組み合わせ（プレゼン標準）
% 和文も明朝にしたい場合は \gtdefault を \mcdefault に変更する
\renewcommand{\kanjifamilydefault}{\gtdefault}
\usefonttheme{serif}

% 行間（1.0 = デフォルト、1.1〜1.2 が「少し広め」の目安）
\linespread{1.15}

% upLaTeX のゴシック体には b（bold）ウェイトが未定義で、
% \bfseries 使用時に警告が出る。b → bx へのマッピングを明示的に
% 宣言することで警告を抑制する（JY2: 横組み、JT2: 縦組み）
\DeclareFontShape{JY2}{hgt}{b}{n}{<->ssub*hgt/bx/n}{}
\DeclareFontShape{JT2}{hgt}{b}{n}{<->ssub*hgt/bx/n}{}

% === Beamer テーマ ===========================================
\usetheme{default}
\usecolortheme{default}

% ナビゲーションバーを削除
\setbeamertemplate{navigation symbols}{}

% === カラー（モノクロ） =======================================
\setbeamercolor{normal text}{fg=black, bg=white}
\setbeamercolor{headline}{fg=black, bg=white}
\setbeamercolor{footline}{fg=black, bg=white}
\setbeamercolor{frametitle}{fg=black, bg=white}
\setbeamercolor{block title}{fg=black, bg=white}
\setbeamercolor{block body}{fg=black, bg=white}
\setbeamercolor{itemize item}{fg=black}
\setbeamercolor{itemize subitem}{fg=black}
\setbeamercolor{itemize subsubitem}{fg=black}
\setbeamercolor{enumerate item}{fg=black}
\setbeamercolor{enumerate subitem}{fg=black}
\setbeamercolor{enumerate subsubitem}{fg=black}
\setbeamercolor{alerted text}{fg=black}

% ブロックの枠をなくす
\setbeamertemplate{blocks}[default]

% === 余白 ====================================================
\setbeamersize{text margin left=5mm, text margin right=5mm}

% === メタデータ（各ファイルで上書き可） ======================
\newcommand{\coursename}{TikZ：経済学でよく使う図}
\author[山田 知明]{山田 知明}
\title{TikZ：経済学でよく使う図}
\institute[明治大学]{明治大学 商学部}
\date{2026年4月7日}

% === ヘッダー（奇数: 左上に講義名 / 偶数: 右上に講義者名） ==
% leftskip/rightskip でマージンを beamercolorbox 内に閉じ込める
\setbeamertemplate{headline}{%
  \begin{beamercolorbox}[%
    wd=\paperwidth, ht=4mm, dp=1.5mm,
    leftskip=5mm, rightskip=5mm]{headline}%
    \ifodd\insertframenumber
      {\footnotesize\coursename}\hfill
    \else
      \hfill
      {\footnotesize\insertshortauthor%
        \ifx\insertshortinstitute\@empty\else
          \ (\insertshortinstitute)%
        \fi}%
    \fi
  \end{beamercolorbox}%
  % 線：テキスト領域幅（左右 5mm マージン分を差し引く）
  \hbox{\hspace{5mm}\rule{\dimexpr\paperwidth-10mm\relax}{0.1pt}}%
}

% === フッター（右下にページ番号） ============================
% [plain] フレーム（タイトルページ等）では自動的に非表示になる
\setbeamertemplate{footline}{%
  \hbox{\hspace{5mm}\rule{\dimexpr\paperwidth-10mm\relax}{0.1pt}}%
  \begin{beamercolorbox}[%
    wd=\paperwidth, ht=2.5mm, dp=1.5mm,
    leftskip=5mm, rightskip=5mm]{footline}%
    %\hfill{\footnotesize\insertframenumber}%
    \hfill{{\footnotesize\insertframenumber\,/\,\inserttotalframenumber}}
  \end{beamercolorbox}%
}

% === フレームタイトル =========================================
\setbeamertemplate{frametitle}{%
  \vspace{1mm}%
  {\large\bfseries\insertframetitle}%
  \ifx\insertframesubtitle\@empty\else
    \\[0.5mm]{\small\itshape\insertframesubtitle}%
  \fi
  \vspace{-1mm}%
}

% === 箇条書き記号 =============================================
\setbeamertemplate{itemize item}{$\bullet$}
\setbeamertemplate{itemize subitem}{$\circ$}
\setbeamertemplate{itemize subsubitem}{--}

% === タイトルページ ===========================================
\setbeamertemplate{title page}{%
  \vfill
  \centering
  {\LARGE\bfseries\inserttitle}\\[5mm]
  {\normalsize\insertauthor}\\[2mm]
  {\normalsize\insertinstitute}\\[1mm]
  {\normalsize\insertdate}%
  \vfill
}

% === 参考文献スタイル =========================================
\bibliographystyle{ecta}

%
%  前提: \documentclass[dvipdfmx, ...]{beamer}
% ============================================================

% === 数学・記号 ==============================================
% amsmath / amsfonts は Beamer が内部でロード済みなので不要
% amssymb のみ Beamer が読まないため明示的にロードする
\usepackage{amssymb}

% === 図版・作図 ==============================================
% standalone: \input で standalone クラスのファイルを取り込む際に
%   プリアンブル（\documentclass 〜 \begin{document}）を自動的に無視し、
%   本文（TikZ コード等）だけを展開する。
%   → 各 TikZ ファイルを単体コンパイルしながら Beamer にも \input できる。
\usepackage{standalone}

\usepackage{tikz}
\usetikzlibrary{arrows.meta}               % Stealth 矢印等
\usetikzlibrary{decorations.pathreplacing} % 波括弧（DiD図等）
\usetikzlibrary{shapes.geometric, positioning}          % 楕円ノード等（因果推論図）

% PGFPlots: 関数プロット・データ可視化
\usepackage{pgfplots}
\pgfplotsset{compat=1.18}
\usepgfplotslibrary{groupplots}       % 複数パネル（上下・左右配置）
\usepgfplotslibrary{fillbetween}      % 2曲線間の塗りつぶし

% === TikZ 図共通カラー ========================================
% standalone プリアンブルは \input 時に無視されるため、
% 各 TikZ ファイルで使う色はここで定義する。
\definecolor{cfcol}  {RGB}{232, 232, 232}     % 観測不可（薄灰）
\definecolor{cfgray}  {RGB}{100, 140, 200}    % 反事実（薄い青）
\definecolor{expgreen} {RGB}{200, 230, 200}   % 拡張局面シェード
\definecolor{fillcol}{RGB}{180, 215, 255}     % 政策効果の差分領域
\definecolor{gdpblue}  {RGB}{30,  90,  180}   % GDP 線・AD 曲線
\definecolor{goldcol}  {RGB}{190, 140,  10}
\definecolor{recred}   {RGB}{240, 200, 200}   % 後退局面シェード
\definecolor{srasgreen}{RGB}{30,  140, 60}    % SRAS（短期総供給曲線）
\definecolor{trendred} {RGB}{190, 40,  40}    % トレンド線・LRAS
\definecolor{obscol} {RGB}{210, 228, 255}     % 観測可能（薄青）
\definecolor{watercol}{RGB}{175, 210, 240}

% graphicx は Beamer が内部でロードするため不要
% subfigure は非推奨（TeX Live が警告を出す）。
\usepackage{subfigure}

% === 参考文献 ================================================
\usepackage{natbib}

% === 欧文フォント =============================================
% newpxtext: Palatino 改良版の本文フォント（mathpazo の現代的後継）
% newpxmath: Palatino 系数式フォント（記号セットが mathpazo より充実）
% \usefonttheme{serif} と組み合わせて Beamer 全体をセリフ体に統一
%
% 代替: Times 系にする場合は以下に切り替え
%   \usepackage{newtxtext}
%   \usepackage{newtxmath}
\usepackage{newpxtext}
\usepackage{newpxmath}

% === 日本語 ==================================================
% otf / pxchfon は \documentclass 直後のメインファイルに記述する想定
% （Beamer のクラスオプション dvipdfmx より先に読まれる必要があるため）
%
%   メインファイルの例:
%     \documentclass[dvipdfmx, t, 10pt]{beamer}
%     \usepackage[deluxe]{otf}
%     \usepackage[hiragino-pron]{pxchfon}
%     % ============================================================
%  my_preamble.sty
%  upLaTeX + dvipdfmx + Beamer 用プリアンブル
%
%  使い方（\documentclass の直後に）:
%    % ============================================================
%  my_preamble.sty
%  upLaTeX + dvipdfmx + Beamer 用プリアンブル
%
%  使い方（\documentclass の直後に）:
%    % ============================================================
%  my_preamble.sty
%  upLaTeX + dvipdfmx + Beamer 用プリアンブル
%
%  使い方（\documentclass の直後に）:
%    \input{my_preamble}
%
%  前提: \documentclass[dvipdfmx, ...]{beamer}
% ============================================================

% === 数学・記号 ==============================================
% amsmath / amsfonts は Beamer が内部でロード済みなので不要
% amssymb のみ Beamer が読まないため明示的にロードする
\usepackage{amssymb}

% === 図版・作図 ==============================================
% standalone: \input で standalone クラスのファイルを取り込む際に
%   プリアンブル（\documentclass 〜 \begin{document}）を自動的に無視し、
%   本文（TikZ コード等）だけを展開する。
%   → 各 TikZ ファイルを単体コンパイルしながら Beamer にも \input できる。
\usepackage{standalone}

\usepackage{tikz}
\usetikzlibrary{arrows.meta}               % Stealth 矢印等
\usetikzlibrary{decorations.pathreplacing} % 波括弧（DiD図等）
\usetikzlibrary{shapes.geometric, positioning}          % 楕円ノード等（因果推論図）

% PGFPlots: 関数プロット・データ可視化
\usepackage{pgfplots}
\pgfplotsset{compat=1.18}
\usepgfplotslibrary{groupplots}       % 複数パネル（上下・左右配置）
\usepgfplotslibrary{fillbetween}      % 2曲線間の塗りつぶし

% === TikZ 図共通カラー ========================================
% standalone プリアンブルは \input 時に無視されるため、
% 各 TikZ ファイルで使う色はここで定義する。
\definecolor{cfcol}  {RGB}{232, 232, 232}     % 観測不可（薄灰）
\definecolor{cfgray}  {RGB}{100, 140, 200}    % 反事実（薄い青）
\definecolor{expgreen} {RGB}{200, 230, 200}   % 拡張局面シェード
\definecolor{fillcol}{RGB}{180, 215, 255}     % 政策効果の差分領域
\definecolor{gdpblue}  {RGB}{30,  90,  180}   % GDP 線・AD 曲線
\definecolor{goldcol}  {RGB}{190, 140,  10}
\definecolor{recred}   {RGB}{240, 200, 200}   % 後退局面シェード
\definecolor{srasgreen}{RGB}{30,  140, 60}    % SRAS（短期総供給曲線）
\definecolor{trendred} {RGB}{190, 40,  40}    % トレンド線・LRAS
\definecolor{obscol} {RGB}{210, 228, 255}     % 観測可能（薄青）
\definecolor{watercol}{RGB}{175, 210, 240}

% graphicx は Beamer が内部でロードするため不要
% subfigure は非推奨（TeX Live が警告を出す）。
\usepackage{subfigure}

% === 参考文献 ================================================
\usepackage{natbib}

% === 欧文フォント =============================================
% newpxtext: Palatino 改良版の本文フォント（mathpazo の現代的後継）
% newpxmath: Palatino 系数式フォント（記号セットが mathpazo より充実）
% \usefonttheme{serif} と組み合わせて Beamer 全体をセリフ体に統一
%
% 代替: Times 系にする場合は以下に切り替え
%   \usepackage{newtxtext}
%   \usepackage{newtxmath}
\usepackage{newpxtext}
\usepackage{newpxmath}

% === 日本語 ==================================================
% otf / pxchfon は \documentclass 直後のメインファイルに記述する想定
% （Beamer のクラスオプション dvipdfmx より先に読まれる必要があるため）
%
%   メインファイルの例:
%     \documentclass[dvipdfmx, t, 10pt]{beamer}
%     \usepackage[deluxe]{otf}
%     \usepackage[hiragino-pron]{pxchfon}
%     \input{my_preamble}

% hyperref の日本語対応（dvipdfmx 環境でのしおり文字化け防止）
\usepackage{pxjahyper}

% === フォントテーマ ==========================================
% 和文: \kanjifamilydefault を \gtdefault にしているのでゴシック体
% 欧文: \usefonttheme{serif} で Palatino（mathpazo）
% → 和文ゴシック + 欧文セリフ の組み合わせ（プレゼン標準）
% 和文も明朝にしたい場合は \gtdefault を \mcdefault に変更する
\renewcommand{\kanjifamilydefault}{\gtdefault}
\usefonttheme{serif}

% 行間（1.0 = デフォルト、1.1〜1.2 が「少し広め」の目安）
\linespread{1.15}

% upLaTeX のゴシック体には b（bold）ウェイトが未定義で、
% \bfseries 使用時に警告が出る。b → bx へのマッピングを明示的に
% 宣言することで警告を抑制する（JY2: 横組み、JT2: 縦組み）
\DeclareFontShape{JY2}{hgt}{b}{n}{<->ssub*hgt/bx/n}{}
\DeclareFontShape{JT2}{hgt}{b}{n}{<->ssub*hgt/bx/n}{}

% === Beamer テーマ ===========================================
\usetheme{default}
\usecolortheme{default}

% ナビゲーションバーを削除
\setbeamertemplate{navigation symbols}{}

% === カラー（モノクロ） =======================================
\setbeamercolor{normal text}{fg=black, bg=white}
\setbeamercolor{headline}{fg=black, bg=white}
\setbeamercolor{footline}{fg=black, bg=white}
\setbeamercolor{frametitle}{fg=black, bg=white}
\setbeamercolor{block title}{fg=black, bg=white}
\setbeamercolor{block body}{fg=black, bg=white}
\setbeamercolor{itemize item}{fg=black}
\setbeamercolor{itemize subitem}{fg=black}
\setbeamercolor{itemize subsubitem}{fg=black}
\setbeamercolor{enumerate item}{fg=black}
\setbeamercolor{enumerate subitem}{fg=black}
\setbeamercolor{enumerate subsubitem}{fg=black}
\setbeamercolor{alerted text}{fg=black}

% ブロックの枠をなくす
\setbeamertemplate{blocks}[default]

% === 余白 ====================================================
\setbeamersize{text margin left=5mm, text margin right=5mm}

% === メタデータ（各ファイルで上書き可） ======================
\newcommand{\coursename}{TikZ：経済学でよく使う図}
\author[山田 知明]{山田 知明}
\title{TikZ：経済学でよく使う図}
\institute[明治大学]{明治大学 商学部}
\date{2026年4月7日}

% === ヘッダー（奇数: 左上に講義名 / 偶数: 右上に講義者名） ==
% leftskip/rightskip でマージンを beamercolorbox 内に閉じ込める
\setbeamertemplate{headline}{%
  \begin{beamercolorbox}[%
    wd=\paperwidth, ht=4mm, dp=1.5mm,
    leftskip=5mm, rightskip=5mm]{headline}%
    \ifodd\insertframenumber
      {\footnotesize\coursename}\hfill
    \else
      \hfill
      {\footnotesize\insertshortauthor%
        \ifx\insertshortinstitute\@empty\else
          \ (\insertshortinstitute)%
        \fi}%
    \fi
  \end{beamercolorbox}%
  % 線：テキスト領域幅（左右 5mm マージン分を差し引く）
  \hbox{\hspace{5mm}\rule{\dimexpr\paperwidth-10mm\relax}{0.1pt}}%
}

% === フッター（右下にページ番号） ============================
% [plain] フレーム（タイトルページ等）では自動的に非表示になる
\setbeamertemplate{footline}{%
  \hbox{\hspace{5mm}\rule{\dimexpr\paperwidth-10mm\relax}{0.1pt}}%
  \begin{beamercolorbox}[%
    wd=\paperwidth, ht=2.5mm, dp=1.5mm,
    leftskip=5mm, rightskip=5mm]{footline}%
    %\hfill{\footnotesize\insertframenumber}%
    \hfill{{\footnotesize\insertframenumber\,/\,\inserttotalframenumber}}
  \end{beamercolorbox}%
}

% === フレームタイトル =========================================
\setbeamertemplate{frametitle}{%
  \vspace{1mm}%
  {\large\bfseries\insertframetitle}%
  \ifx\insertframesubtitle\@empty\else
    \\[0.5mm]{\small\itshape\insertframesubtitle}%
  \fi
  \vspace{-1mm}%
}

% === 箇条書き記号 =============================================
\setbeamertemplate{itemize item}{$\bullet$}
\setbeamertemplate{itemize subitem}{$\circ$}
\setbeamertemplate{itemize subsubitem}{--}

% === タイトルページ ===========================================
\setbeamertemplate{title page}{%
  \vfill
  \centering
  {\LARGE\bfseries\inserttitle}\\[5mm]
  {\normalsize\insertauthor}\\[2mm]
  {\normalsize\insertinstitute}\\[1mm]
  {\normalsize\insertdate}%
  \vfill
}

% === 参考文献スタイル =========================================
\bibliographystyle{ecta}

%
%  前提: \documentclass[dvipdfmx, ...]{beamer}
% ============================================================

% === 数学・記号 ==============================================
% amsmath / amsfonts は Beamer が内部でロード済みなので不要
% amssymb のみ Beamer が読まないため明示的にロードする
\usepackage{amssymb}

% === 図版・作図 ==============================================
% standalone: \input で standalone クラスのファイルを取り込む際に
%   プリアンブル（\documentclass 〜 \begin{document}）を自動的に無視し、
%   本文（TikZ コード等）だけを展開する。
%   → 各 TikZ ファイルを単体コンパイルしながら Beamer にも \input できる。
\usepackage{standalone}

\usepackage{tikz}
\usetikzlibrary{arrows.meta}               % Stealth 矢印等
\usetikzlibrary{decorations.pathreplacing} % 波括弧（DiD図等）
\usetikzlibrary{shapes.geometric, positioning}          % 楕円ノード等（因果推論図）

% PGFPlots: 関数プロット・データ可視化
\usepackage{pgfplots}
\pgfplotsset{compat=1.18}
\usepgfplotslibrary{groupplots}       % 複数パネル（上下・左右配置）
\usepgfplotslibrary{fillbetween}      % 2曲線間の塗りつぶし

% === TikZ 図共通カラー ========================================
% standalone プリアンブルは \input 時に無視されるため、
% 各 TikZ ファイルで使う色はここで定義する。
\definecolor{cfcol}  {RGB}{232, 232, 232}     % 観測不可（薄灰）
\definecolor{cfgray}  {RGB}{100, 140, 200}    % 反事実（薄い青）
\definecolor{expgreen} {RGB}{200, 230, 200}   % 拡張局面シェード
\definecolor{fillcol}{RGB}{180, 215, 255}     % 政策効果の差分領域
\definecolor{gdpblue}  {RGB}{30,  90,  180}   % GDP 線・AD 曲線
\definecolor{goldcol}  {RGB}{190, 140,  10}
\definecolor{recred}   {RGB}{240, 200, 200}   % 後退局面シェード
\definecolor{srasgreen}{RGB}{30,  140, 60}    % SRAS（短期総供給曲線）
\definecolor{trendred} {RGB}{190, 40,  40}    % トレンド線・LRAS
\definecolor{obscol} {RGB}{210, 228, 255}     % 観測可能（薄青）
\definecolor{watercol}{RGB}{175, 210, 240}

% graphicx は Beamer が内部でロードするため不要
% subfigure は非推奨（TeX Live が警告を出す）。
\usepackage{subfigure}

% === 参考文献 ================================================
\usepackage{natbib}

% === 欧文フォント =============================================
% newpxtext: Palatino 改良版の本文フォント（mathpazo の現代的後継）
% newpxmath: Palatino 系数式フォント（記号セットが mathpazo より充実）
% \usefonttheme{serif} と組み合わせて Beamer 全体をセリフ体に統一
%
% 代替: Times 系にする場合は以下に切り替え
%   \usepackage{newtxtext}
%   \usepackage{newtxmath}
\usepackage{newpxtext}
\usepackage{newpxmath}

% === 日本語 ==================================================
% otf / pxchfon は \documentclass 直後のメインファイルに記述する想定
% （Beamer のクラスオプション dvipdfmx より先に読まれる必要があるため）
%
%   メインファイルの例:
%     \documentclass[dvipdfmx, t, 10pt]{beamer}
%     \usepackage[deluxe]{otf}
%     \usepackage[hiragino-pron]{pxchfon}
%     % ============================================================
%  my_preamble.sty
%  upLaTeX + dvipdfmx + Beamer 用プリアンブル
%
%  使い方（\documentclass の直後に）:
%    \input{my_preamble}
%
%  前提: \documentclass[dvipdfmx, ...]{beamer}
% ============================================================

% === 数学・記号 ==============================================
% amsmath / amsfonts は Beamer が内部でロード済みなので不要
% amssymb のみ Beamer が読まないため明示的にロードする
\usepackage{amssymb}

% === 図版・作図 ==============================================
% standalone: \input で standalone クラスのファイルを取り込む際に
%   プリアンブル（\documentclass 〜 \begin{document}）を自動的に無視し、
%   本文（TikZ コード等）だけを展開する。
%   → 各 TikZ ファイルを単体コンパイルしながら Beamer にも \input できる。
\usepackage{standalone}

\usepackage{tikz}
\usetikzlibrary{arrows.meta}               % Stealth 矢印等
\usetikzlibrary{decorations.pathreplacing} % 波括弧（DiD図等）
\usetikzlibrary{shapes.geometric, positioning}          % 楕円ノード等（因果推論図）

% PGFPlots: 関数プロット・データ可視化
\usepackage{pgfplots}
\pgfplotsset{compat=1.18}
\usepgfplotslibrary{groupplots}       % 複数パネル（上下・左右配置）
\usepgfplotslibrary{fillbetween}      % 2曲線間の塗りつぶし

% === TikZ 図共通カラー ========================================
% standalone プリアンブルは \input 時に無視されるため、
% 各 TikZ ファイルで使う色はここで定義する。
\definecolor{cfcol}  {RGB}{232, 232, 232}     % 観測不可（薄灰）
\definecolor{cfgray}  {RGB}{100, 140, 200}    % 反事実（薄い青）
\definecolor{expgreen} {RGB}{200, 230, 200}   % 拡張局面シェード
\definecolor{fillcol}{RGB}{180, 215, 255}     % 政策効果の差分領域
\definecolor{gdpblue}  {RGB}{30,  90,  180}   % GDP 線・AD 曲線
\definecolor{goldcol}  {RGB}{190, 140,  10}
\definecolor{recred}   {RGB}{240, 200, 200}   % 後退局面シェード
\definecolor{srasgreen}{RGB}{30,  140, 60}    % SRAS（短期総供給曲線）
\definecolor{trendred} {RGB}{190, 40,  40}    % トレンド線・LRAS
\definecolor{obscol} {RGB}{210, 228, 255}     % 観測可能（薄青）
\definecolor{watercol}{RGB}{175, 210, 240}

% graphicx は Beamer が内部でロードするため不要
% subfigure は非推奨（TeX Live が警告を出す）。
\usepackage{subfigure}

% === 参考文献 ================================================
\usepackage{natbib}

% === 欧文フォント =============================================
% newpxtext: Palatino 改良版の本文フォント（mathpazo の現代的後継）
% newpxmath: Palatino 系数式フォント（記号セットが mathpazo より充実）
% \usefonttheme{serif} と組み合わせて Beamer 全体をセリフ体に統一
%
% 代替: Times 系にする場合は以下に切り替え
%   \usepackage{newtxtext}
%   \usepackage{newtxmath}
\usepackage{newpxtext}
\usepackage{newpxmath}

% === 日本語 ==================================================
% otf / pxchfon は \documentclass 直後のメインファイルに記述する想定
% （Beamer のクラスオプション dvipdfmx より先に読まれる必要があるため）
%
%   メインファイルの例:
%     \documentclass[dvipdfmx, t, 10pt]{beamer}
%     \usepackage[deluxe]{otf}
%     \usepackage[hiragino-pron]{pxchfon}
%     \input{my_preamble}

% hyperref の日本語対応（dvipdfmx 環境でのしおり文字化け防止）
\usepackage{pxjahyper}

% === フォントテーマ ==========================================
% 和文: \kanjifamilydefault を \gtdefault にしているのでゴシック体
% 欧文: \usefonttheme{serif} で Palatino（mathpazo）
% → 和文ゴシック + 欧文セリフ の組み合わせ（プレゼン標準）
% 和文も明朝にしたい場合は \gtdefault を \mcdefault に変更する
\renewcommand{\kanjifamilydefault}{\gtdefault}
\usefonttheme{serif}

% 行間（1.0 = デフォルト、1.1〜1.2 が「少し広め」の目安）
\linespread{1.15}

% upLaTeX のゴシック体には b（bold）ウェイトが未定義で、
% \bfseries 使用時に警告が出る。b → bx へのマッピングを明示的に
% 宣言することで警告を抑制する（JY2: 横組み、JT2: 縦組み）
\DeclareFontShape{JY2}{hgt}{b}{n}{<->ssub*hgt/bx/n}{}
\DeclareFontShape{JT2}{hgt}{b}{n}{<->ssub*hgt/bx/n}{}

% === Beamer テーマ ===========================================
\usetheme{default}
\usecolortheme{default}

% ナビゲーションバーを削除
\setbeamertemplate{navigation symbols}{}

% === カラー（モノクロ） =======================================
\setbeamercolor{normal text}{fg=black, bg=white}
\setbeamercolor{headline}{fg=black, bg=white}
\setbeamercolor{footline}{fg=black, bg=white}
\setbeamercolor{frametitle}{fg=black, bg=white}
\setbeamercolor{block title}{fg=black, bg=white}
\setbeamercolor{block body}{fg=black, bg=white}
\setbeamercolor{itemize item}{fg=black}
\setbeamercolor{itemize subitem}{fg=black}
\setbeamercolor{itemize subsubitem}{fg=black}
\setbeamercolor{enumerate item}{fg=black}
\setbeamercolor{enumerate subitem}{fg=black}
\setbeamercolor{enumerate subsubitem}{fg=black}
\setbeamercolor{alerted text}{fg=black}

% ブロックの枠をなくす
\setbeamertemplate{blocks}[default]

% === 余白 ====================================================
\setbeamersize{text margin left=5mm, text margin right=5mm}

% === メタデータ（各ファイルで上書き可） ======================
\newcommand{\coursename}{TikZ：経済学でよく使う図}
\author[山田 知明]{山田 知明}
\title{TikZ：経済学でよく使う図}
\institute[明治大学]{明治大学 商学部}
\date{2026年4月7日}

% === ヘッダー（奇数: 左上に講義名 / 偶数: 右上に講義者名） ==
% leftskip/rightskip でマージンを beamercolorbox 内に閉じ込める
\setbeamertemplate{headline}{%
  \begin{beamercolorbox}[%
    wd=\paperwidth, ht=4mm, dp=1.5mm,
    leftskip=5mm, rightskip=5mm]{headline}%
    \ifodd\insertframenumber
      {\footnotesize\coursename}\hfill
    \else
      \hfill
      {\footnotesize\insertshortauthor%
        \ifx\insertshortinstitute\@empty\else
          \ (\insertshortinstitute)%
        \fi}%
    \fi
  \end{beamercolorbox}%
  % 線：テキスト領域幅（左右 5mm マージン分を差し引く）
  \hbox{\hspace{5mm}\rule{\dimexpr\paperwidth-10mm\relax}{0.1pt}}%
}

% === フッター（右下にページ番号） ============================
% [plain] フレーム（タイトルページ等）では自動的に非表示になる
\setbeamertemplate{footline}{%
  \hbox{\hspace{5mm}\rule{\dimexpr\paperwidth-10mm\relax}{0.1pt}}%
  \begin{beamercolorbox}[%
    wd=\paperwidth, ht=2.5mm, dp=1.5mm,
    leftskip=5mm, rightskip=5mm]{footline}%
    %\hfill{\footnotesize\insertframenumber}%
    \hfill{{\footnotesize\insertframenumber\,/\,\inserttotalframenumber}}
  \end{beamercolorbox}%
}

% === フレームタイトル =========================================
\setbeamertemplate{frametitle}{%
  \vspace{1mm}%
  {\large\bfseries\insertframetitle}%
  \ifx\insertframesubtitle\@empty\else
    \\[0.5mm]{\small\itshape\insertframesubtitle}%
  \fi
  \vspace{-1mm}%
}

% === 箇条書き記号 =============================================
\setbeamertemplate{itemize item}{$\bullet$}
\setbeamertemplate{itemize subitem}{$\circ$}
\setbeamertemplate{itemize subsubitem}{--}

% === タイトルページ ===========================================
\setbeamertemplate{title page}{%
  \vfill
  \centering
  {\LARGE\bfseries\inserttitle}\\[5mm]
  {\normalsize\insertauthor}\\[2mm]
  {\normalsize\insertinstitute}\\[1mm]
  {\normalsize\insertdate}%
  \vfill
}

% === 参考文献スタイル =========================================
\bibliographystyle{ecta}


% hyperref の日本語対応（dvipdfmx 環境でのしおり文字化け防止）
\usepackage{pxjahyper}

% === フォントテーマ ==========================================
% 和文: \kanjifamilydefault を \gtdefault にしているのでゴシック体
% 欧文: \usefonttheme{serif} で Palatino（mathpazo）
% → 和文ゴシック + 欧文セリフ の組み合わせ（プレゼン標準）
% 和文も明朝にしたい場合は \gtdefault を \mcdefault に変更する
\renewcommand{\kanjifamilydefault}{\gtdefault}
\usefonttheme{serif}

% 行間（1.0 = デフォルト、1.1〜1.2 が「少し広め」の目安）
\linespread{1.15}

% upLaTeX のゴシック体には b（bold）ウェイトが未定義で、
% \bfseries 使用時に警告が出る。b → bx へのマッピングを明示的に
% 宣言することで警告を抑制する（JY2: 横組み、JT2: 縦組み）
\DeclareFontShape{JY2}{hgt}{b}{n}{<->ssub*hgt/bx/n}{}
\DeclareFontShape{JT2}{hgt}{b}{n}{<->ssub*hgt/bx/n}{}

% === Beamer テーマ ===========================================
\usetheme{default}
\usecolortheme{default}

% ナビゲーションバーを削除
\setbeamertemplate{navigation symbols}{}

% === カラー（モノクロ） =======================================
\setbeamercolor{normal text}{fg=black, bg=white}
\setbeamercolor{headline}{fg=black, bg=white}
\setbeamercolor{footline}{fg=black, bg=white}
\setbeamercolor{frametitle}{fg=black, bg=white}
\setbeamercolor{block title}{fg=black, bg=white}
\setbeamercolor{block body}{fg=black, bg=white}
\setbeamercolor{itemize item}{fg=black}
\setbeamercolor{itemize subitem}{fg=black}
\setbeamercolor{itemize subsubitem}{fg=black}
\setbeamercolor{enumerate item}{fg=black}
\setbeamercolor{enumerate subitem}{fg=black}
\setbeamercolor{enumerate subsubitem}{fg=black}
\setbeamercolor{alerted text}{fg=black}

% ブロックの枠をなくす
\setbeamertemplate{blocks}[default]

% === 余白 ====================================================
\setbeamersize{text margin left=5mm, text margin right=5mm}

% === メタデータ（各ファイルで上書き可） ======================
\newcommand{\coursename}{TikZ：経済学でよく使う図}
\author[山田 知明]{山田 知明}
\title{TikZ：経済学でよく使う図}
\institute[明治大学]{明治大学 商学部}
\date{2026年4月7日}

% === ヘッダー（奇数: 左上に講義名 / 偶数: 右上に講義者名） ==
% leftskip/rightskip でマージンを beamercolorbox 内に閉じ込める
\setbeamertemplate{headline}{%
  \begin{beamercolorbox}[%
    wd=\paperwidth, ht=4mm, dp=1.5mm,
    leftskip=5mm, rightskip=5mm]{headline}%
    \ifodd\insertframenumber
      {\footnotesize\coursename}\hfill
    \else
      \hfill
      {\footnotesize\insertshortauthor%
        \ifx\insertshortinstitute\@empty\else
          \ (\insertshortinstitute)%
        \fi}%
    \fi
  \end{beamercolorbox}%
  % 線：テキスト領域幅（左右 5mm マージン分を差し引く）
  \hbox{\hspace{5mm}\rule{\dimexpr\paperwidth-10mm\relax}{0.1pt}}%
}

% === フッター（右下にページ番号） ============================
% [plain] フレーム（タイトルページ等）では自動的に非表示になる
\setbeamertemplate{footline}{%
  \hbox{\hspace{5mm}\rule{\dimexpr\paperwidth-10mm\relax}{0.1pt}}%
  \begin{beamercolorbox}[%
    wd=\paperwidth, ht=2.5mm, dp=1.5mm,
    leftskip=5mm, rightskip=5mm]{footline}%
    %\hfill{\footnotesize\insertframenumber}%
    \hfill{{\footnotesize\insertframenumber\,/\,\inserttotalframenumber}}
  \end{beamercolorbox}%
}

% === フレームタイトル =========================================
\setbeamertemplate{frametitle}{%
  \vspace{1mm}%
  {\large\bfseries\insertframetitle}%
  \ifx\insertframesubtitle\@empty\else
    \\[0.5mm]{\small\itshape\insertframesubtitle}%
  \fi
  \vspace{-1mm}%
}

% === 箇条書き記号 =============================================
\setbeamertemplate{itemize item}{$\bullet$}
\setbeamertemplate{itemize subitem}{$\circ$}
\setbeamertemplate{itemize subsubitem}{--}

% === タイトルページ ===========================================
\setbeamertemplate{title page}{%
  \vfill
  \centering
  {\LARGE\bfseries\inserttitle}\\[5mm]
  {\normalsize\insertauthor}\\[2mm]
  {\normalsize\insertinstitute}\\[1mm]
  {\normalsize\insertdate}%
  \vfill
}

% === 参考文献スタイル =========================================
\bibliographystyle{ecta}

%
%  前提: \documentclass[dvipdfmx, ...]{beamer}
% ============================================================

% === 数学・記号 ==============================================
% amsmath / amsfonts は Beamer が内部でロード済みなので不要
% amssymb のみ Beamer が読まないため明示的にロードする
\usepackage{amssymb}

% === 図版・作図 ==============================================
% standalone: \input で standalone クラスのファイルを取り込む際に
%   プリアンブル（\documentclass 〜 \begin{document}）を自動的に無視し、
%   本文（TikZ コード等）だけを展開する。
%   → 各 TikZ ファイルを単体コンパイルしながら Beamer にも \input できる。
\usepackage{standalone}

\usepackage{tikz}
\usetikzlibrary{arrows.meta}               % Stealth 矢印等
\usetikzlibrary{decorations.pathreplacing} % 波括弧（DiD図等）
\usetikzlibrary{shapes.geometric, positioning}          % 楕円ノード等（因果推論図）

% PGFPlots: 関数プロット・データ可視化
\usepackage{pgfplots}
\pgfplotsset{compat=1.18}
\usepgfplotslibrary{groupplots}       % 複数パネル（上下・左右配置）
\usepgfplotslibrary{fillbetween}      % 2曲線間の塗りつぶし

% === TikZ 図共通カラー ========================================
% standalone プリアンブルは \input 時に無視されるため、
% 各 TikZ ファイルで使う色はここで定義する。
\definecolor{cfcol}  {RGB}{232, 232, 232}     % 観測不可（薄灰）
\definecolor{cfgray}  {RGB}{100, 140, 200}    % 反事実（薄い青）
\definecolor{expgreen} {RGB}{200, 230, 200}   % 拡張局面シェード
\definecolor{fillcol}{RGB}{180, 215, 255}     % 政策効果の差分領域
\definecolor{gdpblue}  {RGB}{30,  90,  180}   % GDP 線・AD 曲線
\definecolor{goldcol}  {RGB}{190, 140,  10}
\definecolor{recred}   {RGB}{240, 200, 200}   % 後退局面シェード
\definecolor{srasgreen}{RGB}{30,  140, 60}    % SRAS（短期総供給曲線）
\definecolor{trendred} {RGB}{190, 40,  40}    % トレンド線・LRAS
\definecolor{obscol} {RGB}{210, 228, 255}     % 観測可能（薄青）
\definecolor{watercol}{RGB}{175, 210, 240}

% graphicx は Beamer が内部でロードするため不要
% subfigure は非推奨（TeX Live が警告を出す）。
\usepackage{subfigure}

% === 参考文献 ================================================
\usepackage{natbib}

% === 欧文フォント =============================================
% newpxtext: Palatino 改良版の本文フォント（mathpazo の現代的後継）
% newpxmath: Palatino 系数式フォント（記号セットが mathpazo より充実）
% \usefonttheme{serif} と組み合わせて Beamer 全体をセリフ体に統一
%
% 代替: Times 系にする場合は以下に切り替え
%   \usepackage{newtxtext}
%   \usepackage{newtxmath}
\usepackage{newpxtext}
\usepackage{newpxmath}

% === 日本語 ==================================================
% otf / pxchfon は \documentclass 直後のメインファイルに記述する想定
% （Beamer のクラスオプション dvipdfmx より先に読まれる必要があるため）
%
%   メインファイルの例:
%     \documentclass[dvipdfmx, t, 10pt]{beamer}
%     \usepackage[deluxe]{otf}
%     \usepackage[hiragino-pron]{pxchfon}
%     % ============================================================
%  my_preamble.sty
%  upLaTeX + dvipdfmx + Beamer 用プリアンブル
%
%  使い方（\documentclass の直後に）:
%    % ============================================================
%  my_preamble.sty
%  upLaTeX + dvipdfmx + Beamer 用プリアンブル
%
%  使い方（\documentclass の直後に）:
%    \input{my_preamble}
%
%  前提: \documentclass[dvipdfmx, ...]{beamer}
% ============================================================

% === 数学・記号 ==============================================
% amsmath / amsfonts は Beamer が内部でロード済みなので不要
% amssymb のみ Beamer が読まないため明示的にロードする
\usepackage{amssymb}

% === 図版・作図 ==============================================
% standalone: \input で standalone クラスのファイルを取り込む際に
%   プリアンブル（\documentclass 〜 \begin{document}）を自動的に無視し、
%   本文（TikZ コード等）だけを展開する。
%   → 各 TikZ ファイルを単体コンパイルしながら Beamer にも \input できる。
\usepackage{standalone}

\usepackage{tikz}
\usetikzlibrary{arrows.meta}               % Stealth 矢印等
\usetikzlibrary{decorations.pathreplacing} % 波括弧（DiD図等）
\usetikzlibrary{shapes.geometric, positioning}          % 楕円ノード等（因果推論図）

% PGFPlots: 関数プロット・データ可視化
\usepackage{pgfplots}
\pgfplotsset{compat=1.18}
\usepgfplotslibrary{groupplots}       % 複数パネル（上下・左右配置）
\usepgfplotslibrary{fillbetween}      % 2曲線間の塗りつぶし

% === TikZ 図共通カラー ========================================
% standalone プリアンブルは \input 時に無視されるため、
% 各 TikZ ファイルで使う色はここで定義する。
\definecolor{cfcol}  {RGB}{232, 232, 232}     % 観測不可（薄灰）
\definecolor{cfgray}  {RGB}{100, 140, 200}    % 反事実（薄い青）
\definecolor{expgreen} {RGB}{200, 230, 200}   % 拡張局面シェード
\definecolor{fillcol}{RGB}{180, 215, 255}     % 政策効果の差分領域
\definecolor{gdpblue}  {RGB}{30,  90,  180}   % GDP 線・AD 曲線
\definecolor{goldcol}  {RGB}{190, 140,  10}
\definecolor{recred}   {RGB}{240, 200, 200}   % 後退局面シェード
\definecolor{srasgreen}{RGB}{30,  140, 60}    % SRAS（短期総供給曲線）
\definecolor{trendred} {RGB}{190, 40,  40}    % トレンド線・LRAS
\definecolor{obscol} {RGB}{210, 228, 255}     % 観測可能（薄青）
\definecolor{watercol}{RGB}{175, 210, 240}

% graphicx は Beamer が内部でロードするため不要
% subfigure は非推奨（TeX Live が警告を出す）。
\usepackage{subfigure}

% === 参考文献 ================================================
\usepackage{natbib}

% === 欧文フォント =============================================
% newpxtext: Palatino 改良版の本文フォント（mathpazo の現代的後継）
% newpxmath: Palatino 系数式フォント（記号セットが mathpazo より充実）
% \usefonttheme{serif} と組み合わせて Beamer 全体をセリフ体に統一
%
% 代替: Times 系にする場合は以下に切り替え
%   \usepackage{newtxtext}
%   \usepackage{newtxmath}
\usepackage{newpxtext}
\usepackage{newpxmath}

% === 日本語 ==================================================
% otf / pxchfon は \documentclass 直後のメインファイルに記述する想定
% （Beamer のクラスオプション dvipdfmx より先に読まれる必要があるため）
%
%   メインファイルの例:
%     \documentclass[dvipdfmx, t, 10pt]{beamer}
%     \usepackage[deluxe]{otf}
%     \usepackage[hiragino-pron]{pxchfon}
%     \input{my_preamble}

% hyperref の日本語対応（dvipdfmx 環境でのしおり文字化け防止）
\usepackage{pxjahyper}

% === フォントテーマ ==========================================
% 和文: \kanjifamilydefault を \gtdefault にしているのでゴシック体
% 欧文: \usefonttheme{serif} で Palatino（mathpazo）
% → 和文ゴシック + 欧文セリフ の組み合わせ（プレゼン標準）
% 和文も明朝にしたい場合は \gtdefault を \mcdefault に変更する
\renewcommand{\kanjifamilydefault}{\gtdefault}
\usefonttheme{serif}

% 行間（1.0 = デフォルト、1.1〜1.2 が「少し広め」の目安）
\linespread{1.15}

% upLaTeX のゴシック体には b（bold）ウェイトが未定義で、
% \bfseries 使用時に警告が出る。b → bx へのマッピングを明示的に
% 宣言することで警告を抑制する（JY2: 横組み、JT2: 縦組み）
\DeclareFontShape{JY2}{hgt}{b}{n}{<->ssub*hgt/bx/n}{}
\DeclareFontShape{JT2}{hgt}{b}{n}{<->ssub*hgt/bx/n}{}

% === Beamer テーマ ===========================================
\usetheme{default}
\usecolortheme{default}

% ナビゲーションバーを削除
\setbeamertemplate{navigation symbols}{}

% === カラー（モノクロ） =======================================
\setbeamercolor{normal text}{fg=black, bg=white}
\setbeamercolor{headline}{fg=black, bg=white}
\setbeamercolor{footline}{fg=black, bg=white}
\setbeamercolor{frametitle}{fg=black, bg=white}
\setbeamercolor{block title}{fg=black, bg=white}
\setbeamercolor{block body}{fg=black, bg=white}
\setbeamercolor{itemize item}{fg=black}
\setbeamercolor{itemize subitem}{fg=black}
\setbeamercolor{itemize subsubitem}{fg=black}
\setbeamercolor{enumerate item}{fg=black}
\setbeamercolor{enumerate subitem}{fg=black}
\setbeamercolor{enumerate subsubitem}{fg=black}
\setbeamercolor{alerted text}{fg=black}

% ブロックの枠をなくす
\setbeamertemplate{blocks}[default]

% === 余白 ====================================================
\setbeamersize{text margin left=5mm, text margin right=5mm}

% === メタデータ（各ファイルで上書き可） ======================
\newcommand{\coursename}{TikZ：経済学でよく使う図}
\author[山田 知明]{山田 知明}
\title{TikZ：経済学でよく使う図}
\institute[明治大学]{明治大学 商学部}
\date{2026年4月7日}

% === ヘッダー（奇数: 左上に講義名 / 偶数: 右上に講義者名） ==
% leftskip/rightskip でマージンを beamercolorbox 内に閉じ込める
\setbeamertemplate{headline}{%
  \begin{beamercolorbox}[%
    wd=\paperwidth, ht=4mm, dp=1.5mm,
    leftskip=5mm, rightskip=5mm]{headline}%
    \ifodd\insertframenumber
      {\footnotesize\coursename}\hfill
    \else
      \hfill
      {\footnotesize\insertshortauthor%
        \ifx\insertshortinstitute\@empty\else
          \ (\insertshortinstitute)%
        \fi}%
    \fi
  \end{beamercolorbox}%
  % 線：テキスト領域幅（左右 5mm マージン分を差し引く）
  \hbox{\hspace{5mm}\rule{\dimexpr\paperwidth-10mm\relax}{0.1pt}}%
}

% === フッター（右下にページ番号） ============================
% [plain] フレーム（タイトルページ等）では自動的に非表示になる
\setbeamertemplate{footline}{%
  \hbox{\hspace{5mm}\rule{\dimexpr\paperwidth-10mm\relax}{0.1pt}}%
  \begin{beamercolorbox}[%
    wd=\paperwidth, ht=2.5mm, dp=1.5mm,
    leftskip=5mm, rightskip=5mm]{footline}%
    %\hfill{\footnotesize\insertframenumber}%
    \hfill{{\footnotesize\insertframenumber\,/\,\inserttotalframenumber}}
  \end{beamercolorbox}%
}

% === フレームタイトル =========================================
\setbeamertemplate{frametitle}{%
  \vspace{1mm}%
  {\large\bfseries\insertframetitle}%
  \ifx\insertframesubtitle\@empty\else
    \\[0.5mm]{\small\itshape\insertframesubtitle}%
  \fi
  \vspace{-1mm}%
}

% === 箇条書き記号 =============================================
\setbeamertemplate{itemize item}{$\bullet$}
\setbeamertemplate{itemize subitem}{$\circ$}
\setbeamertemplate{itemize subsubitem}{--}

% === タイトルページ ===========================================
\setbeamertemplate{title page}{%
  \vfill
  \centering
  {\LARGE\bfseries\inserttitle}\\[5mm]
  {\normalsize\insertauthor}\\[2mm]
  {\normalsize\insertinstitute}\\[1mm]
  {\normalsize\insertdate}%
  \vfill
}

% === 参考文献スタイル =========================================
\bibliographystyle{ecta}

%
%  前提: \documentclass[dvipdfmx, ...]{beamer}
% ============================================================

% === 数学・記号 ==============================================
% amsmath / amsfonts は Beamer が内部でロード済みなので不要
% amssymb のみ Beamer が読まないため明示的にロードする
\usepackage{amssymb}

% === 図版・作図 ==============================================
% standalone: \input で standalone クラスのファイルを取り込む際に
%   プリアンブル（\documentclass 〜 \begin{document}）を自動的に無視し、
%   本文（TikZ コード等）だけを展開する。
%   → 各 TikZ ファイルを単体コンパイルしながら Beamer にも \input できる。
\usepackage{standalone}

\usepackage{tikz}
\usetikzlibrary{arrows.meta}               % Stealth 矢印等
\usetikzlibrary{decorations.pathreplacing} % 波括弧（DiD図等）
\usetikzlibrary{shapes.geometric, positioning}          % 楕円ノード等（因果推論図）

% PGFPlots: 関数プロット・データ可視化
\usepackage{pgfplots}
\pgfplotsset{compat=1.18}
\usepgfplotslibrary{groupplots}       % 複数パネル（上下・左右配置）
\usepgfplotslibrary{fillbetween}      % 2曲線間の塗りつぶし

% === TikZ 図共通カラー ========================================
% standalone プリアンブルは \input 時に無視されるため、
% 各 TikZ ファイルで使う色はここで定義する。
\definecolor{cfcol}  {RGB}{232, 232, 232}     % 観測不可（薄灰）
\definecolor{cfgray}  {RGB}{100, 140, 200}    % 反事実（薄い青）
\definecolor{expgreen} {RGB}{200, 230, 200}   % 拡張局面シェード
\definecolor{fillcol}{RGB}{180, 215, 255}     % 政策効果の差分領域
\definecolor{gdpblue}  {RGB}{30,  90,  180}   % GDP 線・AD 曲線
\definecolor{goldcol}  {RGB}{190, 140,  10}
\definecolor{recred}   {RGB}{240, 200, 200}   % 後退局面シェード
\definecolor{srasgreen}{RGB}{30,  140, 60}    % SRAS（短期総供給曲線）
\definecolor{trendred} {RGB}{190, 40,  40}    % トレンド線・LRAS
\definecolor{obscol} {RGB}{210, 228, 255}     % 観測可能（薄青）
\definecolor{watercol}{RGB}{175, 210, 240}

% graphicx は Beamer が内部でロードするため不要
% subfigure は非推奨（TeX Live が警告を出す）。
\usepackage{subfigure}

% === 参考文献 ================================================
\usepackage{natbib}

% === 欧文フォント =============================================
% newpxtext: Palatino 改良版の本文フォント（mathpazo の現代的後継）
% newpxmath: Palatino 系数式フォント（記号セットが mathpazo より充実）
% \usefonttheme{serif} と組み合わせて Beamer 全体をセリフ体に統一
%
% 代替: Times 系にする場合は以下に切り替え
%   \usepackage{newtxtext}
%   \usepackage{newtxmath}
\usepackage{newpxtext}
\usepackage{newpxmath}

% === 日本語 ==================================================
% otf / pxchfon は \documentclass 直後のメインファイルに記述する想定
% （Beamer のクラスオプション dvipdfmx より先に読まれる必要があるため）
%
%   メインファイルの例:
%     \documentclass[dvipdfmx, t, 10pt]{beamer}
%     \usepackage[deluxe]{otf}
%     \usepackage[hiragino-pron]{pxchfon}
%     % ============================================================
%  my_preamble.sty
%  upLaTeX + dvipdfmx + Beamer 用プリアンブル
%
%  使い方（\documentclass の直後に）:
%    \input{my_preamble}
%
%  前提: \documentclass[dvipdfmx, ...]{beamer}
% ============================================================

% === 数学・記号 ==============================================
% amsmath / amsfonts は Beamer が内部でロード済みなので不要
% amssymb のみ Beamer が読まないため明示的にロードする
\usepackage{amssymb}

% === 図版・作図 ==============================================
% standalone: \input で standalone クラスのファイルを取り込む際に
%   プリアンブル（\documentclass 〜 \begin{document}）を自動的に無視し、
%   本文（TikZ コード等）だけを展開する。
%   → 各 TikZ ファイルを単体コンパイルしながら Beamer にも \input できる。
\usepackage{standalone}

\usepackage{tikz}
\usetikzlibrary{arrows.meta}               % Stealth 矢印等
\usetikzlibrary{decorations.pathreplacing} % 波括弧（DiD図等）
\usetikzlibrary{shapes.geometric, positioning}          % 楕円ノード等（因果推論図）

% PGFPlots: 関数プロット・データ可視化
\usepackage{pgfplots}
\pgfplotsset{compat=1.18}
\usepgfplotslibrary{groupplots}       % 複数パネル（上下・左右配置）
\usepgfplotslibrary{fillbetween}      % 2曲線間の塗りつぶし

% === TikZ 図共通カラー ========================================
% standalone プリアンブルは \input 時に無視されるため、
% 各 TikZ ファイルで使う色はここで定義する。
\definecolor{cfcol}  {RGB}{232, 232, 232}     % 観測不可（薄灰）
\definecolor{cfgray}  {RGB}{100, 140, 200}    % 反事実（薄い青）
\definecolor{expgreen} {RGB}{200, 230, 200}   % 拡張局面シェード
\definecolor{fillcol}{RGB}{180, 215, 255}     % 政策効果の差分領域
\definecolor{gdpblue}  {RGB}{30,  90,  180}   % GDP 線・AD 曲線
\definecolor{goldcol}  {RGB}{190, 140,  10}
\definecolor{recred}   {RGB}{240, 200, 200}   % 後退局面シェード
\definecolor{srasgreen}{RGB}{30,  140, 60}    % SRAS（短期総供給曲線）
\definecolor{trendred} {RGB}{190, 40,  40}    % トレンド線・LRAS
\definecolor{obscol} {RGB}{210, 228, 255}     % 観測可能（薄青）
\definecolor{watercol}{RGB}{175, 210, 240}

% graphicx は Beamer が内部でロードするため不要
% subfigure は非推奨（TeX Live が警告を出す）。
\usepackage{subfigure}

% === 参考文献 ================================================
\usepackage{natbib}

% === 欧文フォント =============================================
% newpxtext: Palatino 改良版の本文フォント（mathpazo の現代的後継）
% newpxmath: Palatino 系数式フォント（記号セットが mathpazo より充実）
% \usefonttheme{serif} と組み合わせて Beamer 全体をセリフ体に統一
%
% 代替: Times 系にする場合は以下に切り替え
%   \usepackage{newtxtext}
%   \usepackage{newtxmath}
\usepackage{newpxtext}
\usepackage{newpxmath}

% === 日本語 ==================================================
% otf / pxchfon は \documentclass 直後のメインファイルに記述する想定
% （Beamer のクラスオプション dvipdfmx より先に読まれる必要があるため）
%
%   メインファイルの例:
%     \documentclass[dvipdfmx, t, 10pt]{beamer}
%     \usepackage[deluxe]{otf}
%     \usepackage[hiragino-pron]{pxchfon}
%     \input{my_preamble}

% hyperref の日本語対応（dvipdfmx 環境でのしおり文字化け防止）
\usepackage{pxjahyper}

% === フォントテーマ ==========================================
% 和文: \kanjifamilydefault を \gtdefault にしているのでゴシック体
% 欧文: \usefonttheme{serif} で Palatino（mathpazo）
% → 和文ゴシック + 欧文セリフ の組み合わせ（プレゼン標準）
% 和文も明朝にしたい場合は \gtdefault を \mcdefault に変更する
\renewcommand{\kanjifamilydefault}{\gtdefault}
\usefonttheme{serif}

% 行間（1.0 = デフォルト、1.1〜1.2 が「少し広め」の目安）
\linespread{1.15}

% upLaTeX のゴシック体には b（bold）ウェイトが未定義で、
% \bfseries 使用時に警告が出る。b → bx へのマッピングを明示的に
% 宣言することで警告を抑制する（JY2: 横組み、JT2: 縦組み）
\DeclareFontShape{JY2}{hgt}{b}{n}{<->ssub*hgt/bx/n}{}
\DeclareFontShape{JT2}{hgt}{b}{n}{<->ssub*hgt/bx/n}{}

% === Beamer テーマ ===========================================
\usetheme{default}
\usecolortheme{default}

% ナビゲーションバーを削除
\setbeamertemplate{navigation symbols}{}

% === カラー（モノクロ） =======================================
\setbeamercolor{normal text}{fg=black, bg=white}
\setbeamercolor{headline}{fg=black, bg=white}
\setbeamercolor{footline}{fg=black, bg=white}
\setbeamercolor{frametitle}{fg=black, bg=white}
\setbeamercolor{block title}{fg=black, bg=white}
\setbeamercolor{block body}{fg=black, bg=white}
\setbeamercolor{itemize item}{fg=black}
\setbeamercolor{itemize subitem}{fg=black}
\setbeamercolor{itemize subsubitem}{fg=black}
\setbeamercolor{enumerate item}{fg=black}
\setbeamercolor{enumerate subitem}{fg=black}
\setbeamercolor{enumerate subsubitem}{fg=black}
\setbeamercolor{alerted text}{fg=black}

% ブロックの枠をなくす
\setbeamertemplate{blocks}[default]

% === 余白 ====================================================
\setbeamersize{text margin left=5mm, text margin right=5mm}

% === メタデータ（各ファイルで上書き可） ======================
\newcommand{\coursename}{TikZ：経済学でよく使う図}
\author[山田 知明]{山田 知明}
\title{TikZ：経済学でよく使う図}
\institute[明治大学]{明治大学 商学部}
\date{2026年4月7日}

% === ヘッダー（奇数: 左上に講義名 / 偶数: 右上に講義者名） ==
% leftskip/rightskip でマージンを beamercolorbox 内に閉じ込める
\setbeamertemplate{headline}{%
  \begin{beamercolorbox}[%
    wd=\paperwidth, ht=4mm, dp=1.5mm,
    leftskip=5mm, rightskip=5mm]{headline}%
    \ifodd\insertframenumber
      {\footnotesize\coursename}\hfill
    \else
      \hfill
      {\footnotesize\insertshortauthor%
        \ifx\insertshortinstitute\@empty\else
          \ (\insertshortinstitute)%
        \fi}%
    \fi
  \end{beamercolorbox}%
  % 線：テキスト領域幅（左右 5mm マージン分を差し引く）
  \hbox{\hspace{5mm}\rule{\dimexpr\paperwidth-10mm\relax}{0.1pt}}%
}

% === フッター（右下にページ番号） ============================
% [plain] フレーム（タイトルページ等）では自動的に非表示になる
\setbeamertemplate{footline}{%
  \hbox{\hspace{5mm}\rule{\dimexpr\paperwidth-10mm\relax}{0.1pt}}%
  \begin{beamercolorbox}[%
    wd=\paperwidth, ht=2.5mm, dp=1.5mm,
    leftskip=5mm, rightskip=5mm]{footline}%
    %\hfill{\footnotesize\insertframenumber}%
    \hfill{{\footnotesize\insertframenumber\,/\,\inserttotalframenumber}}
  \end{beamercolorbox}%
}

% === フレームタイトル =========================================
\setbeamertemplate{frametitle}{%
  \vspace{1mm}%
  {\large\bfseries\insertframetitle}%
  \ifx\insertframesubtitle\@empty\else
    \\[0.5mm]{\small\itshape\insertframesubtitle}%
  \fi
  \vspace{-1mm}%
}

% === 箇条書き記号 =============================================
\setbeamertemplate{itemize item}{$\bullet$}
\setbeamertemplate{itemize subitem}{$\circ$}
\setbeamertemplate{itemize subsubitem}{--}

% === タイトルページ ===========================================
\setbeamertemplate{title page}{%
  \vfill
  \centering
  {\LARGE\bfseries\inserttitle}\\[5mm]
  {\normalsize\insertauthor}\\[2mm]
  {\normalsize\insertinstitute}\\[1mm]
  {\normalsize\insertdate}%
  \vfill
}

% === 参考文献スタイル =========================================
\bibliographystyle{ecta}


% hyperref の日本語対応（dvipdfmx 環境でのしおり文字化け防止）
\usepackage{pxjahyper}

% === フォントテーマ ==========================================
% 和文: \kanjifamilydefault を \gtdefault にしているのでゴシック体
% 欧文: \usefonttheme{serif} で Palatino（mathpazo）
% → 和文ゴシック + 欧文セリフ の組み合わせ（プレゼン標準）
% 和文も明朝にしたい場合は \gtdefault を \mcdefault に変更する
\renewcommand{\kanjifamilydefault}{\gtdefault}
\usefonttheme{serif}

% 行間（1.0 = デフォルト、1.1〜1.2 が「少し広め」の目安）
\linespread{1.15}

% upLaTeX のゴシック体には b（bold）ウェイトが未定義で、
% \bfseries 使用時に警告が出る。b → bx へのマッピングを明示的に
% 宣言することで警告を抑制する（JY2: 横組み、JT2: 縦組み）
\DeclareFontShape{JY2}{hgt}{b}{n}{<->ssub*hgt/bx/n}{}
\DeclareFontShape{JT2}{hgt}{b}{n}{<->ssub*hgt/bx/n}{}

% === Beamer テーマ ===========================================
\usetheme{default}
\usecolortheme{default}

% ナビゲーションバーを削除
\setbeamertemplate{navigation symbols}{}

% === カラー（モノクロ） =======================================
\setbeamercolor{normal text}{fg=black, bg=white}
\setbeamercolor{headline}{fg=black, bg=white}
\setbeamercolor{footline}{fg=black, bg=white}
\setbeamercolor{frametitle}{fg=black, bg=white}
\setbeamercolor{block title}{fg=black, bg=white}
\setbeamercolor{block body}{fg=black, bg=white}
\setbeamercolor{itemize item}{fg=black}
\setbeamercolor{itemize subitem}{fg=black}
\setbeamercolor{itemize subsubitem}{fg=black}
\setbeamercolor{enumerate item}{fg=black}
\setbeamercolor{enumerate subitem}{fg=black}
\setbeamercolor{enumerate subsubitem}{fg=black}
\setbeamercolor{alerted text}{fg=black}

% ブロックの枠をなくす
\setbeamertemplate{blocks}[default]

% === 余白 ====================================================
\setbeamersize{text margin left=5mm, text margin right=5mm}

% === メタデータ（各ファイルで上書き可） ======================
\newcommand{\coursename}{TikZ：経済学でよく使う図}
\author[山田 知明]{山田 知明}
\title{TikZ：経済学でよく使う図}
\institute[明治大学]{明治大学 商学部}
\date{2026年4月7日}

% === ヘッダー（奇数: 左上に講義名 / 偶数: 右上に講義者名） ==
% leftskip/rightskip でマージンを beamercolorbox 内に閉じ込める
\setbeamertemplate{headline}{%
  \begin{beamercolorbox}[%
    wd=\paperwidth, ht=4mm, dp=1.5mm,
    leftskip=5mm, rightskip=5mm]{headline}%
    \ifodd\insertframenumber
      {\footnotesize\coursename}\hfill
    \else
      \hfill
      {\footnotesize\insertshortauthor%
        \ifx\insertshortinstitute\@empty\else
          \ (\insertshortinstitute)%
        \fi}%
    \fi
  \end{beamercolorbox}%
  % 線：テキスト領域幅（左右 5mm マージン分を差し引く）
  \hbox{\hspace{5mm}\rule{\dimexpr\paperwidth-10mm\relax}{0.1pt}}%
}

% === フッター（右下にページ番号） ============================
% [plain] フレーム（タイトルページ等）では自動的に非表示になる
\setbeamertemplate{footline}{%
  \hbox{\hspace{5mm}\rule{\dimexpr\paperwidth-10mm\relax}{0.1pt}}%
  \begin{beamercolorbox}[%
    wd=\paperwidth, ht=2.5mm, dp=1.5mm,
    leftskip=5mm, rightskip=5mm]{footline}%
    %\hfill{\footnotesize\insertframenumber}%
    \hfill{{\footnotesize\insertframenumber\,/\,\inserttotalframenumber}}
  \end{beamercolorbox}%
}

% === フレームタイトル =========================================
\setbeamertemplate{frametitle}{%
  \vspace{1mm}%
  {\large\bfseries\insertframetitle}%
  \ifx\insertframesubtitle\@empty\else
    \\[0.5mm]{\small\itshape\insertframesubtitle}%
  \fi
  \vspace{-1mm}%
}

% === 箇条書き記号 =============================================
\setbeamertemplate{itemize item}{$\bullet$}
\setbeamertemplate{itemize subitem}{$\circ$}
\setbeamertemplate{itemize subsubitem}{--}

% === タイトルページ ===========================================
\setbeamertemplate{title page}{%
  \vfill
  \centering
  {\LARGE\bfseries\inserttitle}\\[5mm]
  {\normalsize\insertauthor}\\[2mm]
  {\normalsize\insertinstitute}\\[1mm]
  {\normalsize\insertdate}%
  \vfill
}

% === 参考文献スタイル =========================================
\bibliographystyle{ecta}


% hyperref の日本語対応（dvipdfmx 環境でのしおり文字化け防止）
\usepackage{pxjahyper}

% === フォントテーマ ==========================================
% 和文: \kanjifamilydefault を \gtdefault にしているのでゴシック体
% 欧文: \usefonttheme{serif} で Palatino（mathpazo）
% → 和文ゴシック + 欧文セリフ の組み合わせ（プレゼン標準）
% 和文も明朝にしたい場合は \gtdefault を \mcdefault に変更する
\renewcommand{\kanjifamilydefault}{\gtdefault}
\usefonttheme{serif}

% 行間（1.0 = デフォルト、1.1〜1.2 が「少し広め」の目安）
\linespread{1.15}

% upLaTeX のゴシック体には b（bold）ウェイトが未定義で、
% \bfseries 使用時に警告が出る。b → bx へのマッピングを明示的に
% 宣言することで警告を抑制する（JY2: 横組み、JT2: 縦組み）
\DeclareFontShape{JY2}{hgt}{b}{n}{<->ssub*hgt/bx/n}{}
\DeclareFontShape{JT2}{hgt}{b}{n}{<->ssub*hgt/bx/n}{}

% === Beamer テーマ ===========================================
\usetheme{default}
\usecolortheme{default}

% ナビゲーションバーを削除
\setbeamertemplate{navigation symbols}{}

% === カラー（モノクロ） =======================================
\setbeamercolor{normal text}{fg=black, bg=white}
\setbeamercolor{headline}{fg=black, bg=white}
\setbeamercolor{footline}{fg=black, bg=white}
\setbeamercolor{frametitle}{fg=black, bg=white}
\setbeamercolor{block title}{fg=black, bg=white}
\setbeamercolor{block body}{fg=black, bg=white}
\setbeamercolor{itemize item}{fg=black}
\setbeamercolor{itemize subitem}{fg=black}
\setbeamercolor{itemize subsubitem}{fg=black}
\setbeamercolor{enumerate item}{fg=black}
\setbeamercolor{enumerate subitem}{fg=black}
\setbeamercolor{enumerate subsubitem}{fg=black}
\setbeamercolor{alerted text}{fg=black}

% ブロックの枠をなくす
\setbeamertemplate{blocks}[default]

% === 余白 ====================================================
\setbeamersize{text margin left=5mm, text margin right=5mm}

% === メタデータ（各ファイルで上書き可） ======================
\newcommand{\coursename}{TikZ：経済学でよく使う図}
\author[山田 知明]{山田 知明}
\title{TikZ：経済学でよく使う図}
\institute[明治大学]{明治大学 商学部}
\date{2026年4月7日}

% === ヘッダー（奇数: 左上に講義名 / 偶数: 右上に講義者名） ==
% leftskip/rightskip でマージンを beamercolorbox 内に閉じ込める
\setbeamertemplate{headline}{%
  \begin{beamercolorbox}[%
    wd=\paperwidth, ht=4mm, dp=1.5mm,
    leftskip=5mm, rightskip=5mm]{headline}%
    \ifodd\insertframenumber
      {\footnotesize\coursename}\hfill
    \else
      \hfill
      {\footnotesize\insertshortauthor%
        \ifx\insertshortinstitute\@empty\else
          \ (\insertshortinstitute)%
        \fi}%
    \fi
  \end{beamercolorbox}%
  % 線：テキスト領域幅（左右 5mm マージン分を差し引く）
  \hbox{\hspace{5mm}\rule{\dimexpr\paperwidth-10mm\relax}{0.1pt}}%
}

% === フッター（右下にページ番号） ============================
% [plain] フレーム（タイトルページ等）では自動的に非表示になる
\setbeamertemplate{footline}{%
  \hbox{\hspace{5mm}\rule{\dimexpr\paperwidth-10mm\relax}{0.1pt}}%
  \begin{beamercolorbox}[%
    wd=\paperwidth, ht=2.5mm, dp=1.5mm,
    leftskip=5mm, rightskip=5mm]{footline}%
    %\hfill{\footnotesize\insertframenumber}%
    \hfill{{\footnotesize\insertframenumber\,/\,\inserttotalframenumber}}
  \end{beamercolorbox}%
}

% === フレームタイトル =========================================
\setbeamertemplate{frametitle}{%
  \vspace{1mm}%
  {\large\bfseries\insertframetitle}%
  \ifx\insertframesubtitle\@empty\else
    \\[0.5mm]{\small\itshape\insertframesubtitle}%
  \fi
  \vspace{-1mm}%
}

% === 箇条書き記号 =============================================
\setbeamertemplate{itemize item}{$\bullet$}
\setbeamertemplate{itemize subitem}{$\circ$}
\setbeamertemplate{itemize subsubitem}{--}

% === タイトルページ ===========================================
\setbeamertemplate{title page}{%
  \vfill
  \centering
  {\LARGE\bfseries\inserttitle}\\[5mm]
  {\normalsize\insertauthor}\\[2mm]
  {\normalsize\insertinstitute}\\[1mm]
  {\normalsize\insertdate}%
  \vfill
}

% === 参考文献スタイル =========================================
\bibliographystyle{ecta}


% hyperref の日本語対応（dvipdfmx 環境でのしおり文字化け防止）
\usepackage{pxjahyper}

% === フォントテーマ ==========================================
% 和文: \kanjifamilydefault を \gtdefault にしているのでゴシック体
% 欧文: \usefonttheme{serif} で Palatino（mathpazo）
% → 和文ゴシック + 欧文セリフ の組み合わせ（プレゼン標準）
% 和文も明朝にしたい場合は \gtdefault を \mcdefault に変更する
\renewcommand{\kanjifamilydefault}{\gtdefault}
\usefonttheme{serif}

% 行間（1.0 = デフォルト、1.1〜1.2 が「少し広め」の目安）
\linespread{1.15}

% upLaTeX のゴシック体には b（bold）ウェイトが未定義で、
% \bfseries 使用時に警告が出る。b → bx へのマッピングを明示的に
% 宣言することで警告を抑制する（JY2: 横組み、JT2: 縦組み）
\DeclareFontShape{JY2}{hgt}{b}{n}{<->ssub*hgt/bx/n}{}
\DeclareFontShape{JT2}{hgt}{b}{n}{<->ssub*hgt/bx/n}{}

% === Beamer テーマ ===========================================
\usetheme{default}
\usecolortheme{default}

% ナビゲーションバーを削除
\setbeamertemplate{navigation symbols}{}

% === カラー（モノクロ） =======================================
\setbeamercolor{normal text}{fg=black, bg=white}
\setbeamercolor{headline}{fg=black, bg=white}
\setbeamercolor{footline}{fg=black, bg=white}
\setbeamercolor{frametitle}{fg=black, bg=white}
\setbeamercolor{block title}{fg=black, bg=white}
\setbeamercolor{block body}{fg=black, bg=white}
\setbeamercolor{itemize item}{fg=black}
\setbeamercolor{itemize subitem}{fg=black}
\setbeamercolor{itemize subsubitem}{fg=black}
\setbeamercolor{enumerate item}{fg=black}
\setbeamercolor{enumerate subitem}{fg=black}
\setbeamercolor{enumerate subsubitem}{fg=black}
\setbeamercolor{alerted text}{fg=black}

% ブロックの枠をなくす
\setbeamertemplate{blocks}[default]

% === 余白 ====================================================
\setbeamersize{text margin left=5mm, text margin right=5mm}

% === メタデータ（各ファイルで上書き可） ======================
\newcommand{\coursename}{TikZ：経済学でよく使う図}
\author[山田 知明]{山田 知明}
\title{TikZ：経済学でよく使う図}
\institute[明治大学]{明治大学 商学部}
\date{2026年4月7日}

% === ヘッダー（奇数: 左上に講義名 / 偶数: 右上に講義者名） ==
% leftskip/rightskip でマージンを beamercolorbox 内に閉じ込める
\setbeamertemplate{headline}{%
  \begin{beamercolorbox}[%
    wd=\paperwidth, ht=4mm, dp=1.5mm,
    leftskip=5mm, rightskip=5mm]{headline}%
    \ifodd\insertframenumber
      {\footnotesize\coursename}\hfill
    \else
      \hfill
      {\footnotesize\insertshortauthor%
        \ifx\insertshortinstitute\@empty\else
          \ (\insertshortinstitute)%
        \fi}%
    \fi
  \end{beamercolorbox}%
  % 線：テキスト領域幅（左右 5mm マージン分を差し引く）
  \hbox{\hspace{5mm}\rule{\dimexpr\paperwidth-10mm\relax}{0.1pt}}%
}

% === フッター（右下にページ番号） ============================
% [plain] フレーム（タイトルページ等）では自動的に非表示になる
\setbeamertemplate{footline}{%
  \hbox{\hspace{5mm}\rule{\dimexpr\paperwidth-10mm\relax}{0.1pt}}%
  \begin{beamercolorbox}[%
    wd=\paperwidth, ht=2.5mm, dp=1.5mm,
    leftskip=5mm, rightskip=5mm]{footline}%
    %\hfill{\footnotesize\insertframenumber}%
    \hfill{{\footnotesize\insertframenumber\,/\,\inserttotalframenumber}}
  \end{beamercolorbox}%
}

% === フレームタイトル =========================================
\setbeamertemplate{frametitle}{%
  \vspace{1mm}%
  {\large\bfseries\insertframetitle}%
  \ifx\insertframesubtitle\@empty\else
    \\[0.5mm]{\small\itshape\insertframesubtitle}%
  \fi
  \vspace{-1mm}%
}

% === 箇条書き記号 =============================================
\setbeamertemplate{itemize item}{$\bullet$}
\setbeamertemplate{itemize subitem}{$\circ$}
\setbeamertemplate{itemize subsubitem}{--}

% === タイトルページ ===========================================
\setbeamertemplate{title page}{%
  \vfill
  \centering
  {\LARGE\bfseries\inserttitle}\\[5mm]
  {\normalsize\insertauthor}\\[2mm]
  {\normalsize\insertinstitute}\\[1mm]
  {\normalsize\insertdate}%
  \vfill
}

% === 参考文献スタイル =========================================
\bibliographystyle{ecta}
